\documentclass{ieeeaccess}

\input{glyphtounicode}
\pdfgentounicode=1
% PDF 1.4 with a classic cross-reference table (no object streams): the widest
% reader compatibility, for submission systems that parse the file themselves.
\pdfminorversion=4
\pdfobjcompresslevel=0

\usepackage[T1]{fontenc}
\usepackage{cite}
\usepackage{amsmath,amssymb,amsfonts}
\usepackage{booktabs}
\usepackage{graphicx}
\usepackage{textcomp}
\usepackage{url}
\usepackage[hidelinks]{hyperref}

\usepackage{array}
\newcolumntype{R}{>{\raggedleft\arraybackslash}X}

% The class sets a page-wide figure's caption in a one-column tabular, which
% never breaks lines: a caption longer than one line ran 1,157 pt off the page
% and was cut at the margin (measured on Fig. 1). This is the class's
% \@makecaption with that one branch set as a paragraph of the figure's width;
% the table branch and the column-width figure branch are the class's own.
\makeatletter
\long\def\@makecaption#1#2{%
\ifx\@captype\@IEEEtablestring%
\begin{flushleft}
\vspace*{5pt}
{\vss\color{accessblue}\tablecapheadfont #1. \ }{\raggedright\tablecapfont#2\vss}%
\end{flushleft}
\@IEEEtablecaptionsepspace%
\else
\@IEEEfigurecaptionsepspace%
\setbox\@tempboxa\hbox{{\color{accessblue}\figcapheadfont #1. \ }}%
\ifdim \xfigwd >\columnwidth%
\parbox[t]{\linewidth}{\raggedright\noindent\unhbox\@tempboxa\figcapfont#2}%
\else%
{\vss\raggedright\noindent\unhbox\@tempboxa\figcapfont#2\vss}%
\fi\fi\vskip 1pt plus 1pt minus 1pt}
\makeatother

\hypersetup{pdfauthor={},pdfsubject={},pdfcreator={}}
\begin{document}
\history{Date of publication xxxx 00, 0000, date of current version xxxx 00, 0000.}
\doi{10.1109/ACCESS.2017.DOI}

\title{Feasibility Without Preference: What the Quantum Approximate Optimization Algorithm Contributes at a Dominant Penalty Weight}
\author{\uppercase{Anonymous}}

\markboth
{Feasibility Without Preference: What the Quantum Approximate Optimization Algorithm Contributes at a Dominant Penalty Weight}
{Feasibility Without Preference: What the Quantum Approximate Optimization Algorithm Contributes at a Dominant Penalty Weight}




\begin{abstract}
An optimality gap does not show how often a constrained quantum optimizer returns a feasible solution, or which feasible solutions it prefers. We measure both for the depth-2 quantum approximate optimization algorithm (QAOA) on a cardinality-constrained portfolio: 2,565 runs over 30 instances $\times$ 5 initial angles $\times$ 6 penalty weights at 12, 16, and 20 assets, the zero-penalty conditions of a nominal 2,700 merged. The preference tracks one ratio, $\epsilon_{F}$: the objective's spread over the feasible set against the cost's largest Ising coefficient. Under this implementation's dominate-the-objective rule, $\epsilon_{F}$ is 0.015--0.016, where the penalty dominates it and cannot separate equal-cardinality selections. There the optimized state sampled inside that set much as a uniform draw would: median preference ratios 0.998--1.009, median conditional distance 0.001--0.002. Its one-shot success probability stayed below that draw's in all 900 runs at this weight, across the baseline and three control blocks. Optimization raised the feasible probability $P_{F}$ instead. At 20 assets, 1.1 times the feasibility threshold yields $\epsilon_{F}$ 0.98, a median ratio of 3.10, and 41\% of runs above that draw, a share that rises with $\epsilon_{F}$ at every size. Moving the constraint into the ansatz (a Dicke state and a weight-preserving mixer, same depth, steps, and shots) puts 420 of 450 runs above the draw and the exact optimum in 402, against 170 at the dominant weight. Changing the weight to a margin-based one worsens the offset-normalized gap in 29 of 30 instances at 20 assets and raises the 1,000-shot success probability in all 30.
\end{abstract}

\begin{keywords}
benchmark testing, cloud computing, constrained optimization, portfolio optimization, quadratic unconstrained binary optimization (QUBO), quantum approximate optimization algorithm (QAOA), quantum computing, reproducibility of results
\end{keywords}

\titlepgskip=-15pt

\maketitle


\section{Introduction}


Any benchmark for a quantum optimizer has to answer one question: did the quantum part contribute anything to the result? Portfolio selection under a cardinality constraint asks for exactly K of N candidate assets to optimize a risk--return trade-off \cite{ref1}, and it is NP-hard \cite{ref2}. Its usual route to a quantum device is the quadratic unconstrained binary optimization (QUBO), or Ising, encoding \cite{ref3}, \cite{ref4}, surveyed in \cite{ref5}--\cite{ref7}. The usual answer is an optimality gap taken on the energies that the solver returns. That number cannot answer the question. It describes the best sample of one measurement batch, and under a penalty encoding, it is divided by a constant that the formulation chose: one fixed portfolio's gap reads 33.6\% at one weight and 0.22\% at another, while its distance from the optimum in the objective's own range stays at 37.8\% (Sec.~V-D). Of 14 published quantum approximate optimization algorithm (QAOA) and annealing portfolio benchmarks, 10 report no gap, and the four that do use four different normalizations (Supplementary Sec. S-I): the normalizations these benchmarks report are not consistent with each other. That an offset-sensitive score depends on the weight is already known \cite{ref8}--\cite{ref10}.


What the question needs is the probability that a batch of shots contains a feasible portfolio within a quality threshold fixed in advance. A batch is built from one shot: under independent and identically distributed (i.i.d.) shots it holds one with probability 1 $-$ (1 $-$ $P_{\tau}$)$^{S}$. The one-shot success probability is the one that factors, as $P_{\tau}$ = $P_{F}$ $\cdot$ $q_{\tau}$: how often a shot is feasible, times which solutions the sampler prefers among those that are (Sec.~III-C). A uniform draw from the feasible set is the control that makes the second readable: it is what respecting the constraint alone already delivers. We measure both against that control for the QAOA at depth p = 2, on a cross of 30 instances $\times$ 5 initial-angle seeds $\times$ 6 penalty weights at N $\in$ \{12, 16, 20\}, for a total of 2,565 runs, scored against exact optima and exact feasible-set statistics.


One formulation constant moves the second factor: how wide the objective's spread over the feasible set is next to the largest Ising coefficient of the cost, $\epsilon_{F}$ = $\Delta_{F}$ / $s_{H}$ (Sec.~III-A). Under the dominate-the-objective rule that this implementation uses, $\epsilon_{F}$ is 0.015--0.016, and there is little systematic enrichment of the quality set. The optimized state's conditional distribution sits where its initial state already was, and its one-shot success probability remains below that of a uniform feasible draw in all 900 runs of that weight. What optimization changed is the other factor: the feasible probability $P_{F}$ rises on 30 of 30 instance medians at each size. Over the six weights, the share of runs that exceed the control rises monotonically with $\epsilon_{F}$, from 0\% at the dominant weight to 41\% at 1.1 times the feasibility threshold at 20 assets (Sec.~V-C).


That end is also reached by changing the encoding instead of the weight. Carried by the ansatz instead --- the Dicke state $|$$D^{N}_{K}$$\rangle$, a mixer commuting with the Hamming weight \cite{ref11}, \cite{ref12}, and the objective as the whole cost --- the sampler exceeds the control in 420 of 450 runs at the same depth, budget, and shots, and returns the exact optimum in 402 versus 170 at the dominant weight. A second block of 450 runs removes the penalty from the cost but keeps the baseline transverse-field ansatz; because its preference ratios are as large but its feasible probability is not, it remains at 0.4 of the control. Only the product of the two factors rises above the control, and what that control bounds is this penalty encoding, not what a quantum sampler can do. Neither weakening the weight nor changing the ansatz is recommended here, for reasons given in Sec.~V-C.


The reported gap and the acquisition probability are two readings of the same weight, measured on the same runs, and they disagree. Lowering the weight from the dominate-the-objective rule to a margin-based one worsens the gap in 29 of 30 instances at N = 20, and raises, in 30 of 30 instances, the probability that 1,000 shots hold a feasible portfolio within 1\% of the feasible range (+0.086). The gap moves because its denominator moves; the acquisition probability moves with how much of the objective the normalized cost carries, under the settings tested. In that comparison, the gap ranked best the weight at which the optimized state showed no preference.


Scale is a separate question. Sec.~V-G reports the resource counts and the envelope that we measured, and Sec.~VI-A marks which reporting items survive where the feasible set cannot be enumerated. We make no claim about post-classical performance. The question answered here comes before the speedup question: feasibility and preference are what a runtime comparison has to be read against.


\textbf{C1 --- What the circuit contributes, measured at six penalty settings and two ansatzes.} The ingredients are known: scoring on the feasible set, the warning against an oversized penalty, and the Dicke-plus-XY alternative are all documented in \cite{ref8}, \cite{ref11}, \cite{ref12}. What is added is the measurement: both factors, per run, from the exact final state and against a uniform feasible draw, over the six penalty settings of one crossed design and the two A = 0 control blocks beside it (Sec.~V-A, Sec.~V-C).


\textbf{C2 --- What that disagreement costs a benchmark, and what to report instead.} It is quantified on the same runs: 76\% of the runs at 20 assets and the dominant weight pass the gap rule while failing the quality rule (Sec.~V-D). Sec.~VI-A lists what each measurement showed a report has to carry, and which of those items survive where the feasible set cannot be enumerated.


\textbf{C3 --- The procedure and its record.} Every run is tied to its inputs, constants, angles, counts, shots, exact-state scores, and, on metered backends, task counts and dollars, under a manifest of what was planned, what collapsed, what ran, and what failed (Sec.~IV-F). Every aggregation in Sec.~V is re-derivable from the run rows by the script named beside it.


The target of this criticism is the penalty encoding and the reporting conventions, not QAOA. A classical annealer on the same penalty-encoded QUBO reproduces the metric pattern (Sec.~V-D), and the simulated QAOA had no speed advantage at any size that we could run, which agrees with the closest prior benchmark \cite{ref13} (Sec.~V-G). The crossed experiment and its controls are statevector simulations, with four hardware sampling tasks beside them (Sec.~V-F), and the arithmetic was measured instead of assumed \cite{ref14}. At these sizes, it is the draw, not the arithmetic, that decides which portfolio is reported (Sec.~V-F).


\section{Related Work}


\textbf{QUBO portfolio benchmarks.} Markowitz's mean--variance selection \cite{ref1} becomes NP-hard once a cardinality constraint is added \cite{ref2}. The standard route to a quantum device is the QUBO/Ising encoding catalogued by Lucas \cite{ref3}, with the constraint absorbed into a quadratic penalty. Rosenberg et al. \cite{ref15} applied it to multi-period trading on an annealer, and Mugel et al. \cite{ref16} compared annealers, gate-based devices, and tensor networks on 52 assets. Venturelli and Kondratyev \cite{ref17} seed reverse annealing with a greedy solution and note that the cardinality penalty needs a large P that ``is typically associated with precision issues'', which is a hardware-side argument for keeping the penalty scale out of the number being scored. The closest benchmark to this paper is Stopfer and Wagner \cite{ref13}. They evaluate annealing and QAOA against mixed-integer programming and classical heuristics on 250 real-market instances of up to 1,000 assets, restrict the quantum methods to at most 30 assets, and find that classical solvers reach optimality within seconds; our runtime results agree (Sec.~V-G). Uotila et al. \cite{ref18} run 100 DJIA instances against an exact eigensolver. They report QAOA meeting both quality and feasibility criteria on 3 of 100, and they observe spectra agreeing to 10$^{-9}$--10$^{-11}$ while ``the allocations start differing substantially''. Brandhofer et al. \cite{ref8} define the approximation ratio over the feasible set alone, ``independent of the particular method used to enforce the constraint'', and reject the dominate-the-objective penalty rule as ``unnecessarily large''; Sec.~V-D measures what those two corrections are worth rather than rediscovering them. Baker and Radha \cite{ref9} report, in a QAOA portfolio case study, run-to-run fluctuation exceeding finite-shot noise, and they reject the approximation ratio as ``sensitive to \ldots{} the constraint enforcement scheme chosen''. The constant offset that carries that sensitivity is written as c = AK$^{2}$ in \cite{ref10}. How to set the weight is a studied question of its own: Verma and Lewis \cite{ref19} give exact weights that make the unconstrained optima coincide with the feasible ones, Ayodele \cite{ref20} compares static rules, and Montanez-Barrera et al. \cite{ref21} give an unbalanced penalty for inequality constraints that avoids slack variables.


\textbf{What benchmarks report, and what is proposed.} Three documents propose what a benchmark in this field should report, and none of them is a formal standard. Abbas et al. \cite{ref22} prescribe a range-normalized approximation ratio ($C_{\mathrm{max}}$ $-$ $\langle$C$\rangle$)/($C_{\mathrm{max}}$ $-$ $C_{\mathrm{min}}$), with mean-sample and best-sample ratios listed separately, since the best sample ``may improve with a higher number of shots''. The Quantum Optimization Benchmark Library \cite{ref23} mandates run, feasible-run, and successful-run counts against a success threshold, and prescribes scoring on the original problem. It warns that ``the objective values of two QUBOs representing the same problem using different penalty factors cannot be directly compared''. Its smallest portfolio instance has 710 binary variables, which is beyond any statevector simulator. Lozano \cite{ref24} proposes a four-metric protocol, the only one to require billed device time, with a gap regularized by max(1, $|$f*$|$). We adopt the reporting fields of \cite{ref22}, \cite{ref23} where they apply. A census of 14 QAOA and annealing portfolio benchmarks --- \cite{ref8}--\cite{ref10}, \cite{ref13}, \cite{ref15}--\cite{ref18}, \cite{ref25}--\cite{ref30}, selected by a stated protocol (Supplementary Sec. S-I, Table~S1) --- is background for the measurements: 10 of the 14 report no optimality gap, the four that do use four different normalizations, none reports task counts or billed cost, and none of the seven published after \cite{ref22}, \cite{ref23} cites either. What the census counts is that a proposal was not cited. That is not evidence that any paper uses the denominator studied here. Barkoutsos et al. \cite{ref31} aggregate the shot distribution by conditional value-at-risk, because the mean and the minimum are unstable summaries, and best-of-shots is that family's extreme case.


\textbf{Initialization, sampling and the simulation stack.} The classical outer loop is where QAOA is most often limited. Zhou et al. \cite{ref32} introduced the INTERP and FOURIER heuristics, Sack and Serbyn \cite{ref33} showed random initialization to be prone to local minima, and barren plateaus \cite{ref34} together with the NP-hardness of training \cite{ref35} leave no initialization universally safe. Guerreschi and Matsuura \cite{ref36} concluded that hundreds of qubits would be needed before QAOA on Max-Cut could outrun a classical solver, which is the precedent for the runtime result here. Amazon Braket \cite{ref37} exposes simulators and hardware behind one API with per-task and per-shot pricing, and Oralkhan and Zhaxalykov \cite{ref26} use it across trapped-ion and superconducting backends. We run PennyLane \cite{ref38} with the Lightning CPU/GPU statevector devices on NVIDIA cuQuantum \cite{ref39} and adjoint gradients \cite{ref40}. Reproducibility in computer-systems research depends on recording the exact inputs \cite{ref41}--\cite{ref43}, and a cloud quantum experiment adds a billable environment to that list (Sec.~VI).


\textbf{Position of this work.} The formulation, the algorithm and the runtime conclusion are standard. What is added is the measurement: a crossed design (instance $\times$ seed $\times$ penalty), evaluated on both exact final-state probabilities and best-feasible sample quality, with exact references, explicit sampling controls, and a recorded execution environment. We report it because it contradicts what a single-seed, single-instance, single-weight run would have let us claim.


\section{Formulation and Evaluation Quantities}


\subsection{The penalty encoding}


Let x $\in$ \{0,1\}$^{N}$ be a binary selection vector, $\mu$ $\in$ $\mathbb{R}$$^{N}$ the vector of annualized expected returns, and $\Sigma$ $\in$ $\mathbb{R}$$^{N\times N}$ the annualized covariance matrix. With risk-aversion $\lambda$ $\in$ [0,1] and target cardinality K, the cardinality-constrained mean--variance problem


\begin{equation*}
\min\; \lambda\, x^{\mathsf{T}}\Sigma x - (1-\lambda)\, \mu^{\mathsf{T}}x\quad \text{subject to}\quad \sum_i x_i = K
\end{equation*}


is cast as an unconstrained QUBO via a quadratic penalty. We keep four quantities apart, because the paper's central measurement is about which of them a reported number divides by:


\begin{equation*}
f(x) = \lambda\, x^{\mathsf{T}}\hat{\Sigma} x - (1-\lambda)\, \hat{\mu}^{\mathsf{T}}x\quad\text{(the portfolio objective),}
\end{equation*}


\begin{equation*}
V(x) = \Bigl(\sum_i x_i - K\Bigr)^2\quad\text{(the constraint violation),}
\end{equation*}


\begin{equation*}
\begin{aligned} H_A(x) &= f(x) + A\, V(x) \\ &\qquad\text{(the penalized objective the QUBO encodes),}\end{aligned}
\end{equation*}


\begin{equation*}
\begin{aligned} E_A(x) &= H_A(x) - A K^2 \\ &\qquad\text{(what $x^{\mathsf{T}}Qx$ returns once the expansion's}\\ &\qquad\text{constant is dropped).}\end{aligned}
\end{equation*}


Let F be the feasible set, the assignments with $\Sigma$$_i$ x$_i$ = K. On F the penalty vanishes, so $H_{A}$(x) = f(x) and $E_{A}$(x) = f(x) $-$ A$\cdot$K$^{2}$: the $-$A$\cdot$K$^{2}$ is a constant the expansion sheds into an offset, not a penalty paid by feasible solutions. Every energy the results report is $E_{A}$, and every gap denominated by $|$$E_{A}$*$|$ therefore carries that constant. The QAOA circuit sees a fifth quantity, the Ising form of $H_{A}$ rescaled by its largest coefficient (below); we write $\langle$$H_{\mathrm{norm}}$$\rangle$ for its expectation and never interchange it with $H_{A}$ or $E_{A}$, nor compare it between runs at different A.


$\mu$ and $\Sigma$ are normalized by their max-abs value ($\hat{\mu}$, $\hat{\Sigma}$ $\in$ [$-$1,1]) so the penalty A has a dataset-independent meaning. Two consequences are stated now because Sec.~V-D reads a financial metric off the same x. The portfolio evaluated there is the equal-weight one, w = x/K. And because $\mu$ and $\Sigma$ are scaled separately, the relative weight of return and risk in f is not the one $\lambda$ = 0.5 would give on the raw financial units: f is what is optimized, and the Sharpe ratio is measured afterwards on the raw series. And because w = x/K, the objective on a feasible x reads f = $\lambda$K$^{2}$$\cdot$w$^{\mathsf{T}}$$\hat{\Sigma}$w $-$ (1$-$$\lambda$)K$\cdot$$\hat{\mu}$$^{\mathsf{T}}$w, so a fixed $\lambda$ weighs risk against return differently at each K: the sizes are not the same financial risk preference. We use $\lambda$ = 0.5 throughout. Expanding the penalty yields an upper-triangular QUBO matrix Q with


\begin{itemize}
\item diagonal: $Q_{ii}$ = $\lambda$$\cdot$$\hat{\Sigma}_{ii}$ $-$ (1$-$$\lambda$)$\cdot$$\hat{\mu}_{i}$ + A$\cdot$(1 $-$ 2K),
\item off-diagonal: $Q_{ij}$ = 2$\lambda$$\cdot$$\hat{\Sigma}_{ij}$ + 2A,
\end{itemize}


and a constant offset A$\cdot$K$^{2}$. For QAOA the QUBO is mapped to an Ising Hamiltonian by $x_{i}$ = (1 $-$ $z_{i}$)/2, giving linear (h) and quadratic (J) coefficients plus a constant. This layer is device-independent and is unchanged whether the solver is classical or quantum.


\textit{The circuit.} With $s_{H}$ = max($\max_{i}$ $|$$h_{i}$$|$, $\max_{i<j}$ $|$$J_{ij}$$|$) and the normalized cost Hamiltonian


\begin{equation*}
H_C = \frac{1}{s_H}\Bigl(\sum_i h_i Z_i + \sum_{i<j} J_{ij} Z_i Z_j\Bigr),\qquad H_M = -\sum_i X_i,
\end{equation*}


the state at depth p is


\begin{equation*}
|\psi(\gamma,\beta)\rangle = \prod_{\ell=p}^{1} e^{-i\beta_\ell H_M} e^{-i\gamma_\ell H_C}\, H^{\otimes N}|0\rangle^{\otimes N},
\end{equation*}


layers applied cost-first, and the objective minimized is $\langle$$\psi$($\gamma$, $\beta$)$|$ $H_{C}$ $|$$\psi$($\gamma$, $\beta$)$\rangle$ = $\langle$$H_{\mathrm{norm}}$$\rangle$. Terms with $|$coefficient$|$ $\leq$ 10$^{-12}$ are dropped before the circuit is built. Each exponential is one Trotter step of the exact commuting decomposition (every term of $H_{C}$ is diagonal, and $H_{M}$'s terms commute among themselves), giving p$\cdot$[C(N,2) + 2N] parameterized rotations. Angles are drawn once per run as $\gamma_{\ell}$, $\beta_{\ell}$ $\sim$ U(0, $\pi$) from the run's seed and are unbounded during optimization. We use p = 2 for E1--E5 (the depth sweep of Supplementary Sec. S-VI is the exception), so there are four variational parameters. The rescaling by $s_{H}$ leaves argmin unchanged and reported energies use the original Q, but it is not neutral to the search: on the circuit it is the reparameterization $\gamma$ $\rightarrow$ $\gamma$/$s_{H}$, so the same initial-angle range explores a different region of the landscape than the unnormalized problem would.


\subsection{Penalty weights}


Three reference weights recur. The heuristic the formulation uses, and the one we call $A_{\mathrm{heur}}$, is the ``make the penalty dominate the objective'' rule


\begin{equation*}
\begin{aligned} A_{\mathrm{heur}} &= 2 s_{\mathrm{obj}} N + 1,\\ s_{\mathrm{obj}} &= \max\bigl(\lambda \max|\hat{\Sigma}|,\, (1-\lambda)\max|\hat{\mu}|\bigr);\end{aligned}
\end{equation*}


because the normalization makes max$|$$\hat{\Sigma}$$|$ = max$|$$\hat{\mu}$$|$ = 1, at $\lambda$ = 0.5 this is $A_{\mathrm{heur}}$ = N + 1 for every instance in this paper. It is the baseline penalty of this implementation. The \textit{idea} of making the penalty dominate the objective is common; this particular formula is not a field standard and we do not treat it as one. The experiments vary it against reference weights computed from the instance ($A_{\mathrm{crit}}$, $A_{\mathrm{margin}}$), not against competing conventions.


$A_{\mathrm{crit}}$ is the smallest A at which the constrained optimum x* is a global minimizer of $H_{A}$: with $a_{m}$ the minimum of f over assignments of cardinality m,


\begin{equation*}
A_{\mathrm{crit}} = \max\Bigl(0,\; \max_{m \neq K} \frac{f^{*} - a_m}{(m-K)^2}\Bigr).
\end{equation*}


If the inner maximum is positive, an infeasible assignment ties x* at A = $A_{\mathrm{crit}}$ and any larger A separates them strictly; if it is zero, some infeasible assignment already ties x* at A = 0; if it is negative (the clamp is active), x* is strictly better than every infeasible assignment at A = 0 and no tie occurs at any A $\geq$ 0. All 9 of the 30 N = 12 instances with $A_{\mathrm{crit}}$ = 0 are the third case, with inner maxima from $-$9.0$\cdot$10$^{-2}$ to $-$5.9$\cdot$10$^{-4}$, so for them the multiples of $A_{\mathrm{crit}}$ collapse to the single condition A = 0. Ties are decided at \texttt{OPT\_TOL} = 10$^{-9}$ relative to $\Delta_{F}$, not by floating-point equality. For N $\leq$ 20 we compute $a_{m}$ exactly over all 2$^{N}$ assignments; for N $>$ 20 we bound the cardinalities far from K and call the result $A_{\mathrm{safe}}$, an upper bound on $A_{\mathrm{crit}}$.


$A_{\mathrm{margin}}$ follows Brandhofer et al. \cite{ref8} and places the best infeasible assignment at the midpoint between f* and the feasible-set mean $\bar{f}_{F}$,


\begin{equation*}
T = \frac{f^{*} + \bar{f}_F}{2},\quad A_{\mathrm{margin}} = \max\Bigl(0,\; \max_{m \neq K} \frac{T - a_m}{(m-K)^2}\Bigr) + \epsilon_A,
\end{equation*}


with $\epsilon_{A}$ = 10$^{-6}$ recorded per instance. At $A_{\mathrm{margin}}$ every infeasible assignment costs at least T $-$ f* more than x* does, up to $\epsilon_{A}$; the best infeasible assignment sits at T (before $\epsilon_{A}$ is added) when the inner maximum is nonnegative; when it is negative every infeasible assignment already exceeds T at A = 0. Both $A_{\mathrm{crit}}$ and $A_{\mathrm{margin}}$ use exact information about the instance; they are reference conditions, not a procedure we claim scales. Computing them took 0.06 s, 0.8 s and 10.3 s per 30 instances at N = 12, 16 and 20.


\subsection{Evaluation quantities}


All scores are computed by an evaluator that reads x, $\hat{\mu}$, $\hat{\Sigma}$, $\lambda$, K and A, and never reads Q. On every one of the 2$^{N}$ states at N = 8 and N = 12 it agrees with the QUBO path on $H_{A}$, $E_{A}$, the Ising coefficients and the bit order (E0, Sec.~IV-C). With f* = $\min_{F}$ f, $f_{F}$$^{max}$ = $\max_{F}$ f and the feasible-range spread $\Delta_{F}$ = $f_{F}$$^{max}$ $-$ f*, the quantities are those of Table~\ref{tab:I}.


Every score in Sec.~V is one of these, computed by the evaluator from x, $\hat{\mu}$, $\hat{\Sigma}$, $\lambda$, K and A.


\begin{table*}[!t]
\centering
\scriptsize
\caption{The evaluation quantities.}
\label{tab:I}
\begin{tabular}{p{0.141\textwidth}p{0.335\textwidth}p{0.434\textwidth}}
\toprule
Quantity & Definition & Role \\
\midrule
Offset-normalized gap $g_{\mathrm{off}}$ & (f(x) $-$ f*) / $|$f* $-$ A$\cdot$K$^{2}$$|$, feasible x & the gap taken on the solver's own energies; analyzed, not used for judgment \\
Feasible-range gap $g_{F}$ & (f(x) $-$ f*) / $\Delta_{F}$ & primary quality of a feasible solution; free of the constant and of A \\
Feasible probability $P_{F}$ & Pr[x $\in$ F] under the sampler's exact distribution & the chance one shot is usable \\
Quality-hit probability $P_{\tau}$ & Pr[x $\in$ F and $g_{F}$(x) $\leq$ $\tau$] & the unconditional one-shot success probability \\
S-shot success $Q_{\tau}$(S) & 1 $-$ (1 $-$ $P_{\tau}$)$^{S}$ for i.i.d. shots & success at a fixed measurement budget \\
Rank & 1 + number of feasible solutions with f strictly below f(x) & position in F; implies neither a small objective difference nor a similar portfolio \\
Deflation D(A) & $|$f* $-$ A$\cdot$K$^{2}$$|$ / $\Delta_{F}$ = $g_{F}$ / $g_{\mathrm{off}}$ & how many times smaller the offset-normalized number reads \\
Conditional distance $D_{\mathrm{cond}}$ & $\tfrac{1}{2}$ $\sum_{x\in F}$ $|$ p(x)/$P_{F}$ $-$ 1/$|$F$|$ $|$ & how far the sampler, conditioned on feasibility, is from uniform on F \\
Domain metric & return, volatility and Sharpe ratio of w = x/K on the raw series & auxiliary; not what the QUBO optimizes \\
\bottomrule
\end{tabular}
\end{table*}


The primary threshold is $\tau$ = 0.01 of the feasible range; the exact optimum and $\tau$ = 0.05 are sensitivity settings. $\tau$ = 1\% of $\Delta_{F}$ is not the 0.1\% of $|$E*$|$ that the earlier tables used as a success line. The two are compared as decision rules in Sec.~V-D rather than treated as equivalent. Three denominators can vanish, and the evaluator returns the affected quantity as undefined rather than as a number: $\Delta_{F}$ = 0 for $g_{F}$ (the within-$\tau$ test then counts the whole feasible set as qualifying), $|$f* $-$ A$\cdot$K$^{2}$$|$ = 0 for $g_{\mathrm{off}}$ and D(A) together, and $P_{F}$ = 0 for $D_{\mathrm{cond}}$. No setting in Table~\ref{tab:II} reaches any of them. The identity D(A) = $g_{F}$ / $g_{\mathrm{off}}$ holds for non-optimal x; at the optimum both gaps are 0. Runs in which any of these occurs are counted rather than dropped.


Under i.i.d. shots $Q_{\tau}$(S) is a formula in $P_{\tau}$, so shot budgets are compared without re-running QAOA; on hardware, whose shots may be correlated, it is a model. Two consequences follow. It is monotone in $P_{\tau}$ at every S, so one sampler exceeding another at several budgets is one comparison reported at several operating points, decided on $P_{\tau}$. And $Q_{\tau}$(S) is the probability that a batch contains at least one qualifying feasible solution; under the output rule used here it is also the probability of \textit{returning} one. The rule selects the lowest-$E_{A}$ shot among the feasible shots (ties to the lowest basis index), a batch with no feasible shot is a counted outcome rather than a missing value, and $g_{F}$ is not defined for an infeasible x. Because $E_{A}$ and f induce the same ordering on F, the returned shot has $g_{F}$ $\leq$ $\tau$ exactly when the batch holds a feasible shot with $g_{F}$ $\leq$ $\tau$. An infeasible shot of lower $E_{A}$ cannot displace it, so the equivalence survives the weak penalty weights at which the batch minimum is infeasible. It need not hold for a rule that selects the minimum-energy shot without first checking feasibility (\texttt{tests/test\_run\_e1\_cross.py} pins each case). Within one instance at one A, $g_{\mathrm{off}}$ and $g_{F}$ are positive multiples of each other and preserve the order of solutions; what D(A) changes is every comparison across A, across N and against a threshold.


Two identities relate the conditional quantities. Write $u_{\tau}$ = $|$$F_{\tau}$$|$ / $|$F$|$ for the share of the feasible set within $\tau$, and $q_{\tau}$ = $P_{\tau}$ / $P_{F}$ for the sampler's conditional hit rate. Then $|$$q_{\tau}$ $-$ $u_{\tau}$$|$ $\leq$ $D_{\mathrm{cond}}$, and the sampler's one-shot success relative to uniform feasible sampling factors as $P_{\tau}$ / $P_{\tau}$$^{uniform-F}$ = $P_{F}$ $\cdot$ $R_{\tau}$ with $R_{\tau}$ = $q_{\tau}$ / $u_{\tau}$. The first factor is the chance of landing in F, the second the preference for good solutions within it. Neither quantity replaces the comparison. $D_{\mathrm{cond}}$ has no direction: it does not say whether mass moved toward or away from $F_{\tau}$. A small $D_{\mathrm{cond}}$ does not guarantee that $R_{\tau}$ is close to one when $u_{\tau}$ is small ($|$$R_{\tau}$ $-$ 1$|$ $\leq$ $D_{\mathrm{cond}}$ / $u_{\tau}$). Comparisons against uniform-F are therefore made on $P_{\tau}$ directly, with $D_{\mathrm{cond}}$ and $R_{\tau}$ read as their explanation.


\section{Experimental Protocol}


\subsection{Data and problem instances}


Nikkei 225 constituents via \texttt{yfinance}, a three-year window pinned in code to 2023-08-06 -- 2026-08-05 (731 trading days), annualized returns and covariance. For each N the target cardinality is K = max(2, $\lfloor$N/5$\rfloor$). The universe is a fixed list of 50 Nikkei 225 constituents, chosen by the author for liquidity and availability on \texttt{yfinance} rather than by a numerical ranking rule, and written into the source code in the order listed there before any experiment was run. All 50 have a complete daily series in both windows (731 and 735 trading days), so no ticker was excluded for missing data. Instances are subsets of that universe, drawn from a single cached price snapshot. Instance 0 is the \textit{reference instance}, the subset that the single-instance tables use: the first N tickers in the listed order. Instance i $\geq$ 1 is an N-subset drawn without replacement from the 50 by a generator seeded with i, one stream per index, reused at every size. Instances with the same index at different N are therefore not independent draws, since they share 4--7 tickers where independent draws would share 3.8--6.4 in expectation. Instances with different indices are independent draws of subsets from the fixed snapshot, which is not the same as independent market observations. There is no exclusion between instances and no de-duplication. Every instance's sorted ticker list, and a digest of it, are recorded with its results (\texttt{paper01\_e1\_reference.csv}). All E0--E4 instances share one market period and one asset class, so the spread across them measures subset sensitivity rather than robustness across market periods; Supplementary Table~S19 repeats the main contrast without instance 0. E5 repeats a pre-fixed block on the adjacent, non-overlapping window 2020-08-05 -- 2023-08-05 (735 trading days), from a second cached snapshot of the same 50 tickers.


\subsection{Solvers and backends}


\begin{itemize}
\item \textbf{Classical reference.} Exact enumeration of the \textit{feasible} set for N $\leq$ 30 (C(N, K) candidates, 5.9$\cdot$10$^{5}$ instead of 1.07$\cdot$10$^{9}$ at N = 30, so a proven optimum in 5.7 s); the reference solver uses simulated annealing \cite{ref44} only at N = 50 (the annealing arms of Sec.~V-D are controls, not the reference). Every gap reported for N $\leq$ 30 is measured against a proven optimum.
\item \textbf{QAOA baseline.} Depth p = 2, ADAM \cite{ref45} with step 0.1 for 50 steps from uniform random initial angles in [0, $\pi$) drawn from a named seed, gradients by adjoint differentiation \cite{ref40}, 1,000 final measurement shots, on \texttt{lightning.gpu} at single precision. It is the \textit{baseline} of E3 and of the single-seed measurements in the Supplementary Material, not a claim about the best QAOA.
\item \textbf{Simulation host.} RTX A6000 (48 GB VRAM, driver 570.133), Python 3.11, PennyLane 0.45.1 \cite{ref38}, pennylane-lightning-gpu 0.45, custatevec-cu12 1.14 \cite{ref39}. Every number in Sec.~V comes from this host unless a table says otherwise, and numbers from a second machine are never mixed into a table. The same deterministic classical solver's Sharpe ratio differs by 3.4$\cdot$10$^{-16}$ between the Linux host and an Apple-silicon laptop, a BLAS difference.
\item \textbf{Cloud and hardware backends.} The identical code, selected only by a backend name, on AWS Braket SV1 (managed statevector simulator, double precision, us-east-1) and, for the final sampling task only, on IonQ Forte-1 (the optimization loop never runs on hardware).
\end{itemize}


\subsection{The experiments}


Four factors are varied, and they answer different questions. An \textit{instance} is a ticker subset, so the spread over instances measures problem dependence within one universe and one period; an \textit{initial-angle seed} repeats the same problem from a different initialization; a \textit{penalty setting} changes the weight A of the same problem, which moves both the reported metric and the distribution the sampler draws from; the \textit{second price window} (E5) re-runs the same subsets on a different input period. Thirty instances $\times$ five seeds are thus 30 problems each run five times, not 150 problems, and are aggregated as such (Sec.~IV-E).


\textbf{E0 --- evaluator and re-scoring.} The evaluator of Sec.~III-C was checked against the QUBO path on every one of the 2$^{N}$ states at N = 8 and N = 12: $H_{A}$ = $E_{A}$ + A$\cdot$K$^{2}$ on all states, $H_{A}$ = f on F, the same per-cardinality minima, ties and worst feasible value, the same bit order, and the exact $A_{\mathrm{crit}}$ never above the bounded $A_{\mathrm{safe}}$. Fixed feasible solutions of each reference instance were then re-scored without re-running QAOA. Those solutions are the optimum, the solutions of rank 2, 5, 10, 50, 100 and 1,000, the median-rank solution, the worst feasible solution, and the single-seed QAOA solution of Supplementary Table~S2. Each was scored under $g_{\mathrm{off}}$, $g_{F}$ and the all-states approximation ratio of \cite{ref22}, as only A takes the values 1.1, 2, 5 and 10 $\times$ $A_{\mathrm{crit}}$, $A_{\mathrm{margin}}$ and $A_{\mathrm{heur}}$ (354 rows).


\textbf{E1 --- instance $\times$ seed $\times$ penalty cross.} N $\in$ \{12, 16, 20\}; 30 instances per size; initial-angle seeds 42--46; six penalty settings \{1.1, 2, 5, 10\} $\times$ $A_{\mathrm{crit}}$, $A_{\mathrm{margin}}$, $A_{\mathrm{heur}}$; the baseline solver; 1,000 shots. For every run the initial and final angles, the optimization trajectory, the exact probabilities of every feasible bitstring under the initial and the final state, the per-cardinality mass and the shot counts are retained (Sec.~IV-D). The nominal design is 2,700 runs; the nine N = 12 instances with $A_{\mathrm{crit}}$ = 0 collapse their four multiples into one condition A = 0, so 2,565 were expected and 2,565 completed (Table~\ref{tab:II}). Instance 0 at $A_{\mathrm{heur}}$ reproduces an earlier five-seed sweep and the seed-42 $A_{\mathrm{heur}}$ arm an earlier 90-run instance sweep, to the bitstring (energies within 4$\cdot$10$^{-16}$ relative, the ulp of the two evaluators), so those sweeps are not reported separately.


\textbf{E2 --- what the best shot means.} From the retained states of every E1 run, four samplers are compared at S $\in$ \{10, 100, 1000, 10000\}: the optimized QAOA state; the \textit{unoptimized} state at the same initial angles; the uniform distribution over all 2$^{N}$ bitstrings; and the uniform distribution over F. $P_{F}$, $P_{\tau}$ and $D_{\mathrm{cond}}$ are read off the exact state, and $Q_{\tau}$(S) follows by the formula. The uniform-over-F sampler draws K distinct assets uniformly without optimizing the objective. It is a control that uses the constraint directly, matched to QAOA on the number of candidates generated rather than on computational cost, and exact enumeration is used to score its success probability, not to define it. No optimization is re-run.


\textbf{E3 --- optimizer controls.} A block of 300 runs was fixed before any E3 result was seen (instances 0--9, seeds 42--46, \{$A_{\mathrm{heur}}$, $A_{\mathrm{margin}}$\}, at each N), and re-run under two other settings: \textbf{A}, ADAM with step 0.03 for 100 steps; \textbf{B}, SciPy's L-BFGS-B \cite{ref46}, \cite{ref47} on the analytic expectation with adjoint gradients, stopped at SciPy's own criteria or at 100 gradient evaluations, whichever comes first. Each run pairs with its E1 baseline on (N, instance, A, seed) and the same initial angles, so $\langle$$H_{\mathrm{norm}}$$\rangle$ is comparable within a pair. Neither setting is proposed as a better QAOA; the block asks whether the baseline's results are an artifact of one weak setting.


\textbf{E4-A --- the same angles on four backends.} Twelve circuits were fixed before any E4-A result was seen: N $\in$ \{12, 16, 20\} $\times$ instances \{0, 1\} $\times$ seeds \{42, 43\} at $A_{\mathrm{heur}}$, taking the final angles of the stored E1 runs. They are evaluated without optimization on \texttt{lightning.gpu} at double precision (the reference), \texttt{lightning.gpu} at single precision (the precision of every E1 run), \texttt{lightning.qubit} at double precision on one thread, and SV1 at double precision. Each arm returns the exact probabilities of the output distribution and 10 independent batches of 1,000 shots, each batch a separate circuit execution and, on SV1, a separate task. The probabilities are the full 2$^{N}$ vector locally; on SV1, which refuses the \texttt{Probability} result type at zero shots, they are the feasible-set amplitudes from one analytic task. Locally the sampler is re-seeded per batch through the device's own generator (Sec.~VI-C), and SV1's sampler cannot be seeded. The design answers two questions: whether the output distribution at fixed angles is the same, and what repeated sampling from it returns. It does not answer whether an optimization run on SV1 reaches the same angles.


\textbf{E5 --- the E1 block on a second price window.} The E3 block (300 runs) was re-run with the baseline solver on the window that immediately precedes the paper's, from a second cached snapshot of the same 50 tickers under the same instance-drawing rule; the feasible set, $A_{\mathrm{crit}}$, $A_{\mathrm{margin}}$ and the uniform-sampler baselines are recomputed for the new window, which was fixed before any E5 result was seen. The question is whether the direction of the main contrast, the deflation, the disagreement of the two success rules and the E2 conclusion depend on the market period the instances were drawn from. It is a check under a second input condition, not a backtest.


\textbf{E6 --- the constraint in the ansatz instead of the cost.} Every run above carries the cardinality constraint in the cost, as a quadratic penalty, and lets the circuit put amplitude anywhere in 2$^{N}$. E6 reverses that choice and changes nothing else. The initial state is the Dicke state $|$$D^{N}_{K}$$\rangle$ = C(N, K)$^{-1/2}$ $\sum_{|x| = K}$ $|$x$\rangle$, of equal amplitude on every feasible bitstring. The cost is the objective alone, A = 0, normalized by its own $s_{H}$ as every other block is. The mixer layer is a ring of XX + YY terms, each of them exp(i ($\beta$/2) ($X_{i}$ $X_{j}$ + $Y_{i}$ $Y_{j}$)) and each commuting with the total Hamming weight \cite{ref11}, \cite{ref12}; the pair order, the one Trotter sweep it amounts to and the gate the implementation calls are in Supplementary Sec. S-X. The feasible set is then invariant under the whole ansatz, so $P_{F}$ = 1 holds by construction rather than by optimization, and what the angles can still change is the distribution \textit{within} F --- the quantity Sec.~V-A measures and Sec.~V-C tracks. $\epsilon_{F}$ is not 1 here: with no penalty, $s_{H}$ is the objective's own largest Ising coefficient, and $\Delta_{F}$ / $s_{H}$ is 1.27, 1.67 and 2.26 in the instance median (Table~\ref{tab:III}).


The rest is the baseline: 30 instances per size, five initial-angle seeds, ADAM at 0.1 for 50 steps, 1,000 shots, the same evaluator and the same output rule. ``The same budget'' throughout means those three: depth, optimizer steps and shots. It does not mean the same circuit resources: the Dicke preparation, a known construction \cite{ref48} set here as amplitudes, and the N pair gates of each mixer layer are extra to the penalty ansatz's rotation count (Sec.~V-G). A second block of 450 runs changes only the cost: A = 0 with the baseline transverse-field ansatz, the uniform state over 2$^{N}$ and the transverse mixer exp(i $\beta$ $\sum_{i}$ $X_{i}$). It separates the two halves of the change, since E6 moves the constraint into the ansatz \textit{and} takes the penalty out of the cost at the same time.


\textbf{Earlier single-seed series.} The single-instance benchmark at N = 8--30, the depth sweep and the steps $\times$ p $\times$ seed grid, the SV1 cross-backend, matched-precision and repeat runs, the four IonQ Forte-1 tasks, and the classical annealing arms are single-seed or single-instance measurements, and are read as such. Their tables are in the Supplementary Material, and the main text keeps a paragraph on each.


\subsection{Reproducibility of the inputs, and the record}


All conditions are in code, seeds are named and the window is pinned. Four conditions had to be met, each learned by losing results to it. (i) Prices are pinned by \textit{value}, not by date: \texttt{yfinance} revises adjusted closes between calls, so the raw series are cached at full precision and read back with round-trip parsing. The snapshots are third-party data and are not redistributed. The released package carries what the pipeline actually reads: the annualized $\mu$ and $\Sigma$ of each 50-ticker window, $\approx$54 kB each. Every instance's $\mu$ and $\Sigma$ is a slice of those, bit for bit, so every number reproduces from a clone with no network access. (ii) Every instance slices one snapshot, so no two sizes see different revisions of a price. (iii) Statistics are computed on a C-contiguous float64 copy, because \texttt{DataFrame.cov()} sums in memory-layout order. (iv) \texttt{lightning.qubit} is run at \texttt{OMP\_NUM\_THREADS=1}, without which it drifts by $\sim$10$^{-13}$ between runs; \texttt{lightning.gpu} is bitwise reproducible regardless. Figures are regenerated from the committed CSVs with timestamps suppressed. The failures behind each rule --- including the 10$^{-10}$-level input changes that single-precision QAOA amplified into different solutions --- are in Supplementary Sec. S-IX.


Every run is recorded as one CSV row (initial and final angles, evaluation counts, shot counts, exact-state scores, wall-clock and, on metered backends, task identifiers) and every instance as one row with its reference solution, $A_{\mathrm{crit}}$, $A_{\mathrm{margin}}$ and the uniform-sampler baselines. What E1 retains per run is not the full state vector: 2,565 dense initial and final vectors at these sizes would occupy about 16 GB. It keeps the exact probability of every feasible bitstring under the initial and the final state, the per-cardinality mass, the angles, the trajectory and the shot counts, $\approx$100 MB compressed. The E6 block and its control keep the same record for their 900 runs, $\approx$34 MB. Every E2 quantity, $D_{\mathrm{cond}}$ included, is recomputed from those, and the full vector for any run is reconstructed from its recorded angles and seed in seconds.


\subsection{Statistics}


The unit of generalization is the instance. Initial-angle seeds are repeats within an instance, and shots and offline draws are neither. A paired contrast (treatment $-$ control on the same instance, seed and, for E3, initial angles) is first reduced to one value per instance by the median over seeds. The instance-level values are then bootstrapped by resampling instances (2,000 draws, percentile 95\% interval), with the number of instances on which the treatment was better, tied or worse. Two conditions on what the intervals estimate are stated rather than left to be inferred. Because the 30 subsets share one market period, every interval is conditional on that universe, that window and that drawing rule. And because every instance uses the same five seeds, the instance-level value is a contrast \textit{summarized by the median over that fixed set of initial angles}: resampling instances propagates the uncertainty in which subsets were drawn, not the uncertainty in which angles were. Supplementary Table~S19 shows what that costs. At each of the two larger sizes, one of the five seeds reverses the sign of the main contrast on its own, while dropping any one seed leaves it unchanged, so a claim about initializations drawn from a wider distribution would need a design that samples them. The primary contrast is $A_{\mathrm{margin}}$ against $A_{\mathrm{heur}}$, and the primary acquisition quantity is $Q_{0.01}$(1000). $P_{F}$, $g_{F}$ of the best feasible shot, the rate of exact optima and wall-clock are reported beside it, and 5 $\times$ $A_{\mathrm{crit}}$ against $A_{\mathrm{heur}}$ is a secondary contrast on the instances with $A_{\mathrm{crit}}$ $>$ 0. All other cuts are exploratory. Failed runs, runs without a feasible shot, and runs whose lowest-energy shot is infeasible are counted rather than dropped. A contrast on $g_{F}$ is formed only where both sides have a feasible shot, and the number of such pairs is printed next to the estimate. The design was not sized from a minimum detectable effect. The instance count was fixed by the compute budget of the cross rather than by a power calculation: 30 subsets per size, with five seeds and six penalty settings, is 2,565 runs and 13.9 h on one GPU. The resolution that affords is shown with every contrast as ``better in k of 30'', and five seeds show that a within-instance tail exists, not how probable it is.


\subsection{What ran}


E1 rows: nominal is the full cross; ``expected'' applies the $A_{\mathrm{crit}}$ = 0 collapse (nine N = 12 instances run three conditions rather than six); ``failed'' counts execution failures (a run that raised or wrote no result), not optimizer stopping reasons --- an L-BFGS-B line-search failure is a recorded termination, not a failed run (Supplementary Sec. S-VI). ``No feasible shot'' counts completed runs whose 1,000 shots held no feasible bitstring; ``min-energy shot infeasible'' counts runs whose lowest-energy shot violated the constraint although a feasible shot may exist. Time is wall-clock on the host of Sec.~IV-B (E1 includes 0.06, 0.82 and 10.3 s of reference enumeration per size); metered charges are AWS list price as billed, and a ``---'' in that column is an unmetered local run, not a zero cost. The earlier single-seed series are listed in Supplementary Table~S33.


\begin{table*}[!t]
\centering
\scriptsize
\caption{Experiment counts, execution environments and cost.}
\label{tab:II}
\begin{tabular}{p{0.247\textwidth}p{0.101\textwidth}p{0.041\textwidth}p{0.067\textwidth}p{0.101\textwidth}rrr}
\toprule
Experiment & Runs nominal / expected / done / failed & No feasible shot & Min-energy shot infeasible & Environment (precision) & Time & Metered tasks & Metered charge \\
\midrule
E1, N = 12 / 16 / 20 (K = 2/3/4, $|$F$|$ = 66/560/4,845; 9 of 30 N = 12 instances have $A_{\mathrm{crit}}$ = 0) & 2,700 / 2,565 / 2,565 / 0 & 0 / 56 / 1 & 8 / 124 / 95 & lightning.gpu (single) & 13.9 h & --- & --- \\
E0 re-scoring (354 rows) and E2 (from E1 states) & 0 new runs & --- & --- & CPU (double) & minutes & --- & --- \\
E3, ADAM 0.03 $\times$ 100 & 300 / 300 / 300 / 0 & 7 & 8 & lightning.gpu (single) & 3.2 h & --- & --- \\
E3, L-BFGS-B $\leq$ 100 gradients & 300 / 300 / 300 / 0 & 10 & 12 & lightning.gpu (single) & 2.3 h & --- & --- \\
E5, second window & 300 / 300 / 300 / 0 & 11 & 13 & lightning.gpu (single) & 1.6 h & --- & --- \\
E6, constraint in the ansatz (XY mixer, A = 0) & 450 / 450 / 450 / 0 & 0 & 0 & lightning.gpu (single) & 3.5 h & --- & --- \\
E6 control, A = 0 with the baseline transverse-field ansatz & 450 / 450 / 450 / 0 & 0 & 406 & lightning.gpu (single) & 2.5 h & --- & --- \\
E4-A, local arms (3 arms $\times$ 12 circuits $\times$ 11) & 396 / 396 / 396 / 0 & 46 batches & --- & lightning.gpu / .qubit (double, single) & 3 min & --- & --- \\
E4-A, SV1 (12 circuits $\times$ 11 + 2 pilot) & 134 / 134 / 134 / 0 & 14 batches & --- & Braket SV1 (double) & 21 min & 134 & \$0.50 \\
IonQ Forte-1, 100 shots per task & 4 / 4 / 4 / 0 & 0 & --- & Forte-1 & --- & 4 & \$33.20 \\
Earlier SV1 series (Supplementary Sec. S-VII) & 30 & --- & --- & Braket SV1 (double) & --- & $\approx$1,300 & $\approx$\$4.2 \\
Deployment and smoke test & --- & --- & --- & SV1 & --- & 1 & \$0.004 \\
\bottomrule
\end{tabular}
\end{table*}


The Braket total is \$37.90 against a \$50 budget.


\section{Results}


The sections follow the question rather than the order the experiments were run. Sec.~V-A asks what optimization buys at one weight and Sec.~V-B whether that reaches a sampler which only respects the constraint. Sec.~V-C asks what sets the difference between weights and Sec.~V-D what the reported number does over the same range. Sec.~V-E asks whether another optimizer setting or another price window changes it, Sec.~V-F what a second draw from one distribution returns, and Sec.~V-G what the whole measurement costs in resources. Sec.~V-A to Sec.~V-C and Sec.~V-E are the cross and its controls; Sec.~V-D is re-scoring and enumeration, for which no QAOA is re-run; Sec.~V-F is fixed-angle evidence from four backends and single-seed evidence from metered ones, read as such.


\subsection{What optimization changes at the baseline weight (E1)}


\textit{Within the feasible set, the heuristic-weight sampler is close to uniform.} $P_{\tau}$ factors as $P_{F}$ $\cdot$ $q_{\tau}$ with $q_{\tau}$ = Pr[$g_{F}$ $\leq$ $\tau$ $|$ x $\in$ F]. At $A_{\mathrm{heur}}$ the second factor is the share $u_{\tau}$ = $|$$F_{\tau}$$|$/$|$F$|$ that a uniform draw from F would give, and the ratio $R_{\tau}$ = $q_{\tau}$/$u_{\tau}$ is 0.999, 0.998 and 1.009 at N = 12, 16 and 20, with interquartile ranges of 0.02 or less (Supplementary Table~S17). That ratio, however, only checks the mass on one subset. A distribution such as ($\tfrac{1}{2}$, 0, $\tfrac{1}{2}$, 0) on a four-element F has $R_{\tau}$ = 1 on a two-element $F_{\tau}$ and is not uniform. The retained probabilities of every feasible bitstring were therefore used to compute the conditional distance $D_{\mathrm{cond}}$ of Table~\ref{tab:I} directly, per run (Table~\ref{tab:III}; Supplementary Table~S18). At $A_{\mathrm{heur}}$ it is 0.001--0.002 in the median at every size, and never above 0.038 (N = 12), 0.009 (N = 16) or 0.005 (N = 20) in 450 runs. For comparison, a point mass sits at 0.985--1.000, and a 1,000-shot empirical draw from the uniform distribution itself would show 0.10--0.81 at these $|$F$|$. Small is not zero: the distribution is \textit{near}-uniform, and because $D_{\mathrm{cond}}$ carries no direction, the departures from uniform that it measures are read through the $R_{\tau}$ above. At $A_{\mathrm{margin}}$ the conditional distribution departs further from uniform ($D_{\mathrm{cond}}$ 0.17, 0.14 and 0.09 in the median). Enrichment of the primary quality set is clear at N = 12 and 16 (median $R_{\tau}$ 1.40 and 1.55), whereas at N = 20 the median $R_{\tau}$ stays close to one (1.01) although the distribution is not uniform. The multiples of $A_{\mathrm{crit}}$ depart more the closer A is to $A_{\mathrm{crit}}$ ($D_{\mathrm{cond}}$ 0.37--0.49 at 1.1$\cdot$$A_{\mathrm{crit}}$).


\textit{What optimization does at a fixed weight, and what the weight does.} The two are separated by pairing each run's final state with its initial state (Supplementary Table~S34). At $A_{\mathrm{heur}}$, optimization raises the feasible probability $P_{F}$: $P_{F}$ goes from 0.114, 0.035 and 0.087 in the median initial state to 0.179, 0.229 and 0.167, a paired median increase of +0.070, +0.128 and +0.134 on 30 of 30 instances at N = 12, 16 and 20. The expected constraint violation E[V] falls with it, from 21.9, 34.8 and 45.7 to 10.6, 4.0 and 14.3. The conditional distribution stays where it started: $D_{\mathrm{cond}}$ moves from 0.002--0.004 to 0.001--0.002 in the median, and the initial states that were as far as 0.41 from uniform on F for some seeds end within 0.005 of it in every N = 20 run. At $A_{\mathrm{heur}}$, optimization therefore raised the feasible probability $P_{F}$ with little systematic enrichment of the primary quality set.


Separately, changing the weight from $A_{\mathrm{heur}}$ to $A_{\mathrm{margin}}$ after optimization raises the median $P_{F}$ from 0.167 to 0.303 at N = 20 (Table~\ref{tab:IV}). At $A_{\mathrm{margin}}$, optimization raises $P_{F}$ as well (+0.18, +0.22 and +0.24), and the final conditional distribution departs from uniformity, with the enrichment of Table~\ref{tab:III} at N $\leq$ 16. As a complementary multiplicative summary of the same matched runs, $P_{\tau}$ = $P_{F}$ $\cdot$ $u_{\tau}$ $\cdot$ $R_{\tau}$ on each (instance, seed) pair splits the geometric-mean ratio of $P_{\tau}$ into a feasible-mass factor and a conditional-enrichment factor. The mass accounts for 60\%, 21\% and 70\% of the mean log-ratio at N = 12, 16 and 20, and the $R_{\tau}$ ratio exceeds one on 113, 125 and 88 of the 150 pairs (Supplementary Table~S38). Those shares are of the mean log-ratio, so they complement the median contrasts of Table~\ref{tab:IV} rather than decomposing them, and the geometric-mean $R_{\tau}$ ratio at N = 20 is a different statistic from the median $R_{\tau}$ of Table~\ref{tab:III}.


The paired increases at $A_{\mathrm{heur}}$ are nearly the same on every instance (interquartile range $\leq$ 0.001), which is what a circuit dominated by the permutation-symmetric penalty term would give. With f removed, the cost, the mixer and the initial state are all invariant under qubit permutations, so every bitstring of one cardinality has the same probability at any angles, and the conditional distribution on F is exactly uniform. The measured $\epsilon_{F}$(A) = $\Delta_{F}$ / $s_{H}$ puts the objective's feasible range in units of $s_{H}$, the largest Ising coefficient of the cost before it is divided by $s_{H}$ (Sec.~III-A); equally, $\epsilon_{F}$ is the spread of the normalized cost over F. In the instance median it is 0.015, 0.015 and 0.016 at $A_{\mathrm{heur}}$, with the objective's own largest Ising coefficient at 0.7--1.2\% of $s_{H}$ (Supplementary Table~S35), against 0.45--0.52 at $A_{\mathrm{margin}}$ and 1.27--2.26 at A = 0, where the objective is the whole cost. That scale is consistent with the near-uniformity, and it is offered as context rather than as a demonstration: the symmetry argument is exact only for the penalty alone, and how far a 1.5\% objective perturbation is resolved by a given optimizer at finite A is what $D_{\mathrm{cond}}$ measures. Most of the $A_{\mathrm{heur}}$ state's mass remains far from cardinality K after optimization, and $P_{F}$ is what best-of-shots has to rescue.


\textit{The seed axis, within every instance.} The within-instance spread over five seeds of $Q_{\tau}$(1000) has a median over instances of 0.86, 0.49 and 0.19 at $A_{\mathrm{heur}}$, which is most of the quantity's range at the two smaller sizes. At N = 16 and $A_{\mathrm{heur}}$, 14 of 150 runs returned no feasible bitstring among 1,000 shots (9\%; under 1\% at N = 12 and N = 20), a within-instance tail rather than a size trend. Across settings, the lowest-energy shot is infeasible in 8, 124 and 95 runs (Table~\ref{tab:II}), and in 55\% of N = 20 runs at 1.1$\cdot$$A_{\mathrm{crit}}$. At that setting the best \textit{feasible} shot nevertheless has the smallest $g_{F}$ of any setting ($-$3.3 points of $\Delta_{F}$ against $A_{\mathrm{heur}}$, 29 of 30 instances), so the selection rule of Sec.~III-C is not a formality.


\subsection{Comparison with uniform feasible sampling (E2)}


\textit{What the best shot means (E2; Fig.~\ref{fig:1}).} The same exact states answer the question that the best-of-shots rule begs: how often does a batch of S shots contain a solution within $\tau$, and how does that compare with not optimizing at all? Optimization \textit{works}, in the sense that the literature's convergence curves show. At $A_{\mathrm{heur}}$ the optimized state raises $Q_{\tau}$(1000) over the unoptimized state at the same initial angles, from 0.062 to 0.336 at N = 16 and from 0.019 to 0.059 at N = 20, and it beats the uniform-over-2$^{N}$ sampler by one to two orders of magnitude. Its one-shot success probability at the primary threshold $\tau$ = 1\%, however, never exceeds that of a uniform draw from the feasible set: $P_{\tau}$ $<$ $P_{\tau}$$^{uniform-F}$ in 450 of 450 $A_{\mathrm{heur}}$ runs, and therefore at every S from 10 to 10,000. That is one direct comparison of $P_{\tau}$, shown at four operating points because $Q_{\tau}$(S) is monotone in $P_{\tau}$. It is a measured inequality rather than a consequence of the small $D_{\mathrm{cond}}$ above (Sec.~III-C), and it holds at the two sensitivity thresholds as well, 0 of 450 at the exact optimum and at $\tau$ = 5\% (Supplementary Table~S14).

\Figure[!t](topskip=0pt, botskip=0pt, midskip=0pt)[width=\textwidth]{figures/fig_e2_shots.pdf}{E2. $Q_{\tau}$(S) at $\tau$ = 1\% against the shot budget S $\in$ \{10, 100, 1,000, 10,000\} for four samplers: the optimized QAOA state, the unoptimized state at the same initial angles, uniform over F, and uniform over all bitstrings. Panels are (a) N = 12, (b) N = 16 and (c) N = 20, all at $A_{\mathrm{heur}}$. Each instance is first summarized by the median over its five initial-angle seeds; the line is the median of those instance values over the 30 instances and the band their interquartile range. The two uniform samplers are instance properties evaluated on the same instances, and $Q_{\tau}$(S) follows from the exact one-shot probabilities under i.i.d. shots.\label{fig:1}}



At S = 1,000 and N = 20, a uniform feasible draw succeeds with probability 0.34 and the optimized state with 0.06. A benchmark that reported its best-of-1,000 gap would not show the difference at N = 12, where both samplers usually contain a good shot (median $Q_{\tau}$(1000) 0.93 and 1.00). At N = 16 the gap between them is already wide (0.34 against 0.83), and at N = 20 neither sampler reliably contains one (levels at every S: Supplementary Table~S14). The comparison depends on the penalty setting: at $A_{\mathrm{margin}}$ the optimized state beats the feasible-uniform sampler in 7--9\% of runs, and at 1.1$\cdot$$A_{\mathrm{crit}}$ in 41\% of N = 20 runs, so the 0 of 450 is a result at $A_{\mathrm{heur}}$. At these sizes and budgets, then, a best-of-shots figure alone cannot confirm that the optimization contributed, while the exact-state comparison with explicit control samplers can. What that comparison shows at $A_{\mathrm{heur}}$ and the evaluated thresholds is that optimization improves acquisition relative to the unoptimized state, and that the resulting sampler stays far better than an uninformed one and below uniform feasible sampling. The uniform-over-F sampler is not a computational peer of QAOA (Sec.~IV-C). It is the baseline that a constraint-respecting method has to clear, and the classical annealers of Sec.~V-D clear it by a wide margin.


\textit{Multistart at equal shots.} The five E1 seeds can be re-read as a five-start strategy: optimize five times, draw 1,000 shots from each state, and keep the best feasible candidate of the 5,000. The comparison is with what a user can run at the same shot budget, which is to optimize once from a random seed and sample 5,000 times, with Q computed per seed and then averaged. By the AM--GM inequality the pooled strategy is never worse at equal shots, so the question is the increment against four extra optimizations. It is +0.046 to +0.336 in the instance median over the six cells, above 0.01 on 30 of 30 instances in every cell but one, for 41--118 s of extra optimization per instance (Supplementary Table~S24). Five initial angles do buy acquisition probability that five thousand shots from one start do not. None of the strategies reaches uniform sampling from the feasible set at the same 5,000 candidates, nor the annealing baseline of Sec.~V-D.


\subsection{What sets the preference: the weight, and where the constraint lives (E1, E6)}


The two factors answer for one weight. The cross ran six, and the weight is the one formulation constant that moves them: Supplementary Fig.~S13 shows the four quantities of Sec.~III-C against the penalty setting, one series per size. What the weight sets is how wide the objective's feasible range is next to $s_{H}$, which is the $\epsilon_{F}$ = $\Delta_{F}$ / $s_{H}$ of Sec.~V-A and equally the spread of the normalized cost over F, and the preference inside the feasible set follows it.


The threshold is $\tau$ = 1\% of the feasible range throughout. Per (N, A) over the runs of that cell: 150 runs, 30 instances $\times$ 5 seeds, except the positive multiples of $A_{\mathrm{crit}}$ at N = 12, which have 105 over the 21 instances with $A_{\mathrm{crit}}$ $>$ 0, since $A_{\mathrm{crit}}$ = 0 collapses the multiples to the same A = 0 setting. Rows are ordered by $\epsilon_{F}$ = $\Delta_{F}$ / $s_{H}$, the objective's spread over F divided by $s_{H}$ = max($\max_{i}$ $|$$h_{i}$$|$, $\max_{i<j}$ $|$$J_{ij}$$|$), measured per instance (\texttt{coefficient\_scale.py}) and taken as the median over the instances the cell ran. $\epsilon_{F}$ is a ratio and not a share: it passes 1 wherever the objective's range exceeds that coefficient, as the A = 0 rows do. $P_{F}$ and $R_{\tau}$ are the two factors of $P_{\tau}$ / $P_{\tau}$$^{uniform-F}$ = $P_{F}$ $\cdot$ $R_{\tau}$ (Sec.~III-C), as medians over runs; ``above uniform-F'' counts runs whose one-shot success probability exceeds a uniform draw's from F, the comparison E2 makes at $A_{\mathrm{heur}}$; ``infeasible minima'' counts runs whose lowest-energy shot is infeasible, which the feasibility-aware output rule of Sec.~III-C absorbs. $g_{F}$ is of the best feasible shot of the 1,000; $D_{\mathrm{cond}}$ is from the same runs (Supplementary Tables~S17--S18 give quartiles and the initial state). Every column is a median taken separately over the runs of its cell, so multiplying the printed $P_{F}$ and $R_{\tau}$ need not reproduce the printed ratio. $|$F$|$ is 66, 560 and 4,845. The last six rows are E6 (Sec.~IV-C). In both of them the cost is the objective alone, so their $\epsilon_{F}$ is the A = 0 value, measured on the same 30 instances per size; $D_{\mathrm{cond}}$ is computed from that block's own artifacts. They differ in the ansatz, and therefore in $P_{F}$, which is 1 by construction for the constraint-preserving one and measured for the published one. The N = 12 cells do not all cover the same instances, so the monotonicity there is read on the 21 that ran every weight (\texttt{paper01\_e1\_penalty\_ladder\_common.csv}, Sec.~V-C); at N = 16 and N = 20 the two files agree row for row.


\begin{table*}[!t]
\centering
\scriptsize
\caption{Feasibility, conditional preference and sampling success across penalty settings and ansatzes.}
\label{tab:III}
\begin{tabular}{rlrrrrrrrrr}
\toprule
N & A & $\epsilon_{F}$ & $P_{F}$ & $R_{\tau}$ & $D_{\mathrm{cond}}$ & $P_{\tau}$ $\div$ uniform-F & Above uniform-F & $Q_{\tau}$(1000) & $g_{F}$ & Infeasible minima \\
\midrule
12 & $A_{\mathrm{heur}}$ & 0.015 & 0.179 & 1.00 & 0.002 & 0.17 & 0/150 (0\%) & 0.928 & 0.000 & 0/150 \\
12 & 10$\cdot$$A_{\mathrm{crit}}$ & 0.248 & 0.190 & 1.13 & 0.108 & 0.17 & 3/105 (3\%) & 0.929 & 0.000 & 0/105 \\
12 & 5$\cdot$$A_{\mathrm{crit}}$ & 0.396 & 0.244 & 1.23 & 0.207 & 0.26 & 6/105 (6\%) & 0.992 & 0.000 & 0/105 \\
12 & $A_{\mathrm{margin}}$ & 0.520 & 0.328 & 1.40 & 0.173 & 0.47 & 13/150 (9\%) & 1.000 & 0.000 & 1/150 \\
12 & 2$\cdot$$A_{\mathrm{crit}}$ & 0.663 & 0.277 & 1.40 & 0.308 & 0.46 & 10/105 (10\%) & 0.999 & 0.000 & 3/105 \\
12 & 1.1$\cdot$$A_{\mathrm{crit}}$ & 0.780 & 0.257 & 1.83 & 0.382 & 0.47 & 12/105 (11\%) & 0.999 & 0.000 & 3/105 \\
16 & $A_{\mathrm{heur}}$ & 0.015 & 0.229 & 1.00 & 0.001 & 0.23 & 0/150 (0\%) & 0.336 & 0.017 & 14/150 \\
16 & 10$\cdot$$A_{\mathrm{crit}}$ & 0.184 & 0.309 & 1.10 & 0.080 & 0.36 & 1/150 (1\%) & 0.488 & 0.017 & 15/150 \\
16 & 5$\cdot$$A_{\mathrm{crit}}$ & 0.328 & 0.359 & 1.27 & 0.106 & 0.47 & 6/150 (4\%) & 0.584 & 0.013 & 18/150 \\
16 & $A_{\mathrm{margin}}$ & 0.455 & 0.301 & 1.55 & 0.139 & 0.42 & 11/150 (7\%) & 0.536 & 0.015 & 14/150 \\
16 & 2$\cdot$$A_{\mathrm{crit}}$ & 0.621 & 0.275 & 1.65 & 0.192 & 0.43 & 22/150 (15\%) & 0.553 & 0.003 & 16/150 \\
16 & 1.1$\cdot$$A_{\mathrm{crit}}$ & 0.868 & 0.245 & 2.26 & 0.369 & 0.58 & 33/150 (22\%) & 0.661 & 0.000 & 47/150 \\
20 & $A_{\mathrm{heur}}$ & 0.016 & 0.167 & 1.01 & 0.002 & 0.17 & 0/150 (0\%) & 0.059 & 0.062 & 1/150 \\
20 & 10$\cdot$$A_{\mathrm{crit}}$ & 0.176 & 0.219 & 0.97 & 0.036 & 0.30 & 1/150 (1\%) & 0.083 & 0.049 & 0/150 \\
20 & 5$\cdot$$A_{\mathrm{crit}}$ & 0.325 & 0.233 & 0.97 & 0.061 & 0.36 & 5/150 (3\%) & 0.085 & 0.048 & 1/150 \\
20 & $A_{\mathrm{margin}}$ & 0.473 & 0.303 & 1.01 & 0.090 & 0.43 & 10/150 (7\%) & 0.108 & 0.042 & 0/150 \\
20 & 2$\cdot$$A_{\mathrm{crit}}$ & 0.663 & 0.224 & 1.55 & 0.198 & 0.45 & 37/150 (25\%) & 0.132 & 0.032 & 11/150 \\
20 & 1.1$\cdot$$A_{\mathrm{crit}}$ & 0.976 & 0.238 & 3.10 & 0.494 & 0.80 & 62/150 (41\%) & 0.224 & 0.019 & 82/150 \\
12 & XY, A = 0 & 1.266 & 1.000 & 4.49 & 0.470 & 4.49 & 145/150 (97\%) & 1.000 & 0.000 & 0/150 \\
16 & XY, A = 0 & 1.669 & 1.000 & 4.74 & 0.515 & 4.74 & 134/150 (89\%) & 1.000 & 0.000 & 0/150 \\
20 & XY, A = 0 & 2.258 & 1.000 & 9.89 & 0.535 & 9.89 & 141/150 (94\%) & 0.941 & 0.000 & 0/150 \\
12 & A = 0, transverse & 1.266 & 0.149 & 3.03 & 0.519 & 0.41 & 38/150 (25\%) & 0.999 & 0.000 & 106/150 \\
16 & A = 0, transverse & 1.669 & 0.091 & 4.32 & 0.520 & 0.40 & 31/150 (21\%) & 0.556 & 0.000 & 150/150 \\
20 & A = 0, transverse & 2.258 & 0.058 & 8.45 & 0.539 & 0.45 & 28/150 (19\%) & 0.137 & 0.030 & 150/150 \\
\bottomrule
\end{tabular}
\end{table*}


\textit{When the circuit orders the feasible set, and when it does not.} Asking E2's question at every weight that the cross ran makes the near-uniformity at $A_{\mathrm{heur}}$ one end of a monotone relation rather than an isolated fact (Table~\ref{tab:III} and Fig.~\ref{fig:2}, from \texttt{analyze\_ladder.py}). Ordered by $\epsilon_{F}$, the share of runs whose one-shot success probability exceeds a uniform draw from F rises from 0\% at $A_{\mathrm{heur}}$ to 11\%, 22\% and 41\% at 1.1$\cdot$$A_{\mathrm{crit}}$ (N = 12, 16, 20), monotonically at every size and over six weights. At N = 12 the six cells do not cover the same instances, because $A_{\mathrm{crit}}$ = 0 on 9 of the 30. On the 21 that ran every weight the ordering is the same, and the shares are 0\%, 2.9\%, 5.7\%, 7.6\%, 9.5\% and 11.4\%, with $R_{\tau}$ from 1.00 to 1.83 (\texttt{paper01\_e1\_penalty\_ladder\_common.csv}). At N = 16 and N = 20 that file equals Table~\ref{tab:III} row for row. $R_{\tau}$ rises with the share: 1.00 to 1.83 at N = 12 and 1.00 to 2.26 at N = 16, while at N = 20 it stays within 3\% of one up to $\epsilon_{F}$ = 0.47 and then reaches 3.10 at the weakest weight. The best feasible shot's $g_{F}$ falls overall with the same ordering at N = 16 and 20, though not at every step (Table~\ref{tab:III}). At N = 20, $\epsilon_{F}$ is 0.016 at $A_{\mathrm{heur}}$ and 0.98 at 1.1$\cdot$$A_{\mathrm{crit}}$.

\Figure[!t](topskip=0pt, botskip=0pt, midskip=0pt)[width=\textwidth]{figures/fig_ladder_preference.pdf}{Conditional preference and sampling success against a uniform feasible draw, over the penalty settings and the two A = 0 blocks, at $\tau$ = 1\% (Sec.~V-C). The horizontal coordinate is $\epsilon_{F}$ = $\Delta_{F}$ / $s_{H}$, the objective's feasible range divided by the largest Ising coefficient of the cost before normalization, one series per N. (a) The preference ratio $R_{\tau}$ inside the feasible set, median over the runs of each setting with the interquartile range, and the line $R_{\tau}$ = 1 that a uniform draw would give. (b) The share of runs whose one-shot success probability exceeds a uniform draw from the feasible set. Each series is labeled at its ends, $A_{\mathrm{heur}}$ at the left and 1.1, meaning 1.1$\cdot$$A_{\mathrm{crit}}$, at the right. The stars are the two A = 0 blocks at their measured $\epsilon_{F}$, filled for the constraint-preserving ansatz and hollow for the published one. Both panels draw the cells of Table~\ref{tab:III} as they stand, whose N = 12 cells cover different instances (Sec.~V-C); the A = 0 cell of the cross is left out there, because at a threshold of zero every multiple is the same A = 0 setting.\label{fig:2}}



Three facts keep this from being a recommendation. The median one-shot probability stays below that draw even at the weakest weight (0.80 of it at N = 20), so what rises with $\epsilon_{F}$ is the share of runs that cross it, not the median run. $A_{\mathrm{crit}}$ is computed by enumeration (Sec.~III-B), so a weight set as a multiple of it is a reference condition rather than a procedure that scales. And the weakest weight is where the batch minimum is infeasible most often, in 55\% of runs at N = 20 and 31\% at N = 16, so it is usable only under a rule that selects the best feasible shot rather than the lowest-energy one (Sec.~III-C). What the ladder does establish is that the preference the circuit expresses inside F tracks $\epsilon_{F}$, and that the dominate-the-objective rule puts the weight at the end of that range where the preference is absent and the offset-normalized gap reads at its smallest.


\textit{With the constraint in the ansatz, the circuit clears the baseline.} The ladder varies one formulation constant and keeps the encoding. E6 changes the encoding instead (Sec.~IV-C): the Dicke state, a weight-preserving ring mixer, and the objective as the whole cost. $P_{F}$ = 1 then holds by construction, while $\epsilon_{F}$ is measured at the A = 0 scales like every other row of Table~\ref{tab:III}. It is 1.27, 1.67 and 2.26 in the instance median, above 1 because the objective's feasible range is wider than its own largest Ising coefficient. Over 450 runs on the same instances and seeds at the baseline budget, the one-shot success probability is 4.5, 4.7 and 9.9 times that of the uniform feasible draw in the median (interquartile 3.1--5.8, 2.5--7.5 and 6.0--13.4 at N = 12, 16 and 20), and above that draw in 145, 134 and 141 of 150 runs. The conditional distribution sits 0.47--0.54 from uniform on F, where the dominant weight leaves it at 0.002. The 1,000-shot batch returns the exact optimum in 149, 142 and 111 of 150 runs, against 115, 51 and 4 of 150 at the dominant weight. No batch of the block had an infeasible lowest-energy shot, because no amplitude leaves F at any angle (\texttt{tests/test\_xy\_mixer.py} pins that property rather than the numbers). The uniform feasible draw therefore bounds what this penalty encoding leaves the circuit able to do, not on what a quantum sampler can do on this problem.


The comparison has two halves, and the control separates them. Removing the penalty from the cost but keeping the baseline transverse-field ansatz (A = 0, transverse mixer, 450 runs) buys the \textit{preference}: $R_{\tau}$ is 3.03, 4.32 and 8.45 in the median, within the spread of E6's 4.49, 4.74 and 9.89, and the conditional distribution sits 0.52--0.54 from uniform. What it does not buy is the \textit{mass}. $P_{F}$ falls to 0.149, 0.091 and 0.058, since nothing now pushes the state toward F, so the product stays at 0.41, 0.40 and 0.45 of that draw's and crosses it in 97 of 450 runs. The lowest-energy shot is infeasible in 106 of 150 runs at N = 12 and in every run at N = 16 and 20, which the feasibility-aware output rule absorbs (Sec.~III-C). The two factors of Sec.~III-C therefore moved apart between these two encodings: with the penalty gone, the preference is there and the mass is not, and only the product clears the control. That is not a demonstration that the two can be set independently, since the ladder moves both at once, with $P_{F}$ varying from 0.167 to 0.359 across the penalty settings of Table~\ref{tab:III}. It is a measurement of what each encoding gave.


\subsection{The penalty weight and the reported gap (E0, E1)}


The same runs carry the number a benchmark would report, and it moves over this range in the opposite direction. Two measurements separate what the number does from what the solver does: re-scoring \textit{fixed} feasible solutions as only A changes (E0, Fig.~\ref{fig:3}), and the paired contrasts of the cross at two reference weights (Table~\ref{tab:IV}, by the statistics of Sec.~IV-E).

\Figure[!t](topskip=0pt, botskip=0pt, midskip=0pt)[width=\textwidth]{figures/fig_e0_rescore.pdf}{E0. (a) The median-rank feasible solution of the reference N = 20 instance scored three ways as only A changes: the offset-normalized gap $g_{\mathrm{off}}$ = (f $-$ f*)/$|$f* $-$ AK$^{2}$$|$, the all-states range gap 1 $-$ AR = ($E_{A}$ $-$ $E_{A}$*)/($C_{\mathrm{max}}$ $-$ $E_{A}$*) of \cite{ref22}, where $C_{\mathrm{max}}$ is the largest $E_{A}$ over all 2$^{N}$ assignments, and the feasible-range gap $g_{F}$ = (f $-$ f*)/$\Delta_{F}$. (b) The deflation D(A) = $|$f* $-$ A$\cdot$K$^{2}$$|$/$\Delta_{F}$ per size on the reference instances, with $A_{\mathrm{heur}}$ marked by a star; N $>$ 20 uses $A_{\mathrm{safe}}$ and is drawn hollow. (c) D($A_{\mathrm{heur}}$) over 30 instances per size (box = IQR, whisker = range), with the reference instance marked.\label{fig:3}}



The two level columns and the contrast column answer different questions: levels are medians over the runs of a setting, the contrast is the median of within-instance median paired differences, and the difference of two levels is therefore not the paired contrast. Paired $A_{\mathrm{margin}}$ $-$ $A_{\mathrm{heur}}$ on the same instance and seed, reduced to an instance median, then bootstrapped over instances (2,000 resamples, 95\% percentile interval); ``better'' counts instances on which $A_{\mathrm{margin}}$ is better, with ties in parentheses. Gap differences are in percentage points, and a count of pairs is of runs with a feasible shot on both sides, not of instances. Gaps are of the best \textit{feasible} shot and are formed only where both sides have one (pairs shown); $P_{F}$, $P_{\tau}$ and $Q_{\tau}$(1000) are exact-state quantities and exist for every run.


\begin{table*}[!t]
\centering
\scriptsize
\caption{Main effects of the penalty setting.}
\label{tab:IV}
\begin{tabular}{rp{0.260\textwidth}rrp{0.224\textwidth}p{0.133\textwidth}}
\toprule
N & Quantity & $A_{\mathrm{heur}}$ level & $A_{\mathrm{margin}}$ level & Paired median diff [95\% CI] & Better (tied) of 30 \\
\midrule
12 & $P_{F}$ & 0.179 & 0.328 & +0.006 [+0.005, +0.012] & 28 \\
12 & $P_{\tau}$ & 2.6$\cdot$10$^{-3}$ & 7.6$\cdot$10$^{-3}$ & +2.1$\cdot$10$^{-3}$ [+1.3$\cdot$10$^{-3}$, +3.3$\cdot$10$^{-3}$] & 27 \\
12 & $Q_{\tau}$(1000) & 0.928 & 0.9995 & +0.011 [+0.005, +0.034] & 27 \\
12 & $g_{\mathrm{off}}$ (\%) & 0.000 & 0.000 & 0 [0, 0] & 0 (29) \\
12 & $g_{F}$ (\%) & 0.00 & 0.00 & 0 [0, 0] & 1 (28) \\
16 & $P_{F}$ & 0.229 & 0.301 & +0.0007 [+0.0004, +0.0010] & 27 \\
16 & $P_{\tau}$ & 4.1$\cdot$10$^{-4}$ & 7.7$\cdot$10$^{-4}$ & +8.2$\cdot$10$^{-5}$ [+6.9$\cdot$10$^{-5}$, +1.4$\cdot$10$^{-4}$] & 30 \\
16 & $Q_{\tau}$(1000) & 0.336 & 0.536 & +0.040 [+0.031, +0.061] & 30 \\
16 & $g_{\mathrm{off}}$ (\%), 129 pairs & 0.0143 & 0.5437 & \textbf{+0.44 [+0.18, +0.75]} (worse) & 4 (3) \\
16 & $g_{F}$ (\%), 129 pairs & 1.70 & 1.49 & 0.00 [$-$0.43, 0.00] & 13 (8) \\
20 & $P_{F}$ & 0.167 & 0.303 & +0.034 [+0.017, +0.056] & 30 \\
20 & $P_{\tau}$ & 6.1$\cdot$10$^{-5}$ & 1.1$\cdot$10$^{-4}$ & +9.2$\cdot$10$^{-5}$ [+7.6$\cdot$10$^{-5}$, +1.4$\cdot$10$^{-4}$] & 30 \\
20 & $Q_{\tau}$(1000) & 0.059 & 0.108 & \textbf{+0.086 [+0.071, +0.116]} & 30 \\
20 & $g_{\mathrm{off}}$ (\%), 149 pairs & 0.039 & 0.964 & \textbf{+0.81 [+0.64, +1.11]} (worse) & 1 (0) \\
20 & $g_{F}$ (\%), 149 pairs & 6.23 & 4.22 & \textbf{$-$1.61 [$-$3.11, $-$0.67]} & 24 (1) \\
20 & Rank of best feasible shot, 149 pairs & 23 & 10.5 & $-$11.8 [$-$17.5, $-$3.0] & 24 (1) \\
\bottomrule
\end{tabular}
\end{table*}


\textit{The offset-normalized gap and the sampler move in opposite directions.} Replacing $A_{\mathrm{heur}}$ by $A_{\mathrm{margin}}$ makes $g_{\mathrm{off}}$ \textit{worse} in 29 of 30 instances at N = 20 (median +0.81 percentage points) and 23 of 30 at N = 16. Every acquisition quantity moves the other way: $P_{F}$, $P_{\tau}$ and $Q_{\tau}$(1000) are better in 30 of 30 instances at N = 20 and in 27--30 of 30 at the other sizes. The feasible-range gap of the best feasible shot is better in 24 of 30 at N = 20 (median $-$1.6 points of $\Delta_{F}$) and unchanged in the median at N = 16 (better in 13, tied in 8, worse in 9). The secondary contrast, 5$\cdot$$A_{\mathrm{crit}}$ against $A_{\mathrm{heur}}$ on the instances with $A_{\mathrm{crit}}$ $>$ 0, agrees on $Q_{\tau}$(1000) (better in 29 of 30 at N = 16, +0.094 [+0.069, +0.113]; 30 of 30 at N = 20, +0.019 [+0.002, +0.045]) and on the direction of $g_{\mathrm{off}}$. It is mixed on the best feasible shot's $g_{F}$ (worse in 17 of 30 at N = 16, better in 23 of 30 at N = 20). Which weight yields the best single shot and which yields the highest $Q_{\tau}$(1000) are different questions.


\textit{What moves the gap is the denominator.} The paired ratio of offset-normalized gaps factors, pair by pair, into a solution part (f($x_{{A_{2}}}$) $-$ f*)/(f($x_{{A_{1}}}$) $-$ f*) and a denominator part $|$f* $-$ A$_1$K$^{2}$$|$/$|$f* $-$ A$_2$K$^{2}$$|$. The identity does not survive a median, so the decomposition is taken in the geometric mean over the same pairs. At N = 20, on the 128 pairs with a positive gap on both sides, the $A_{\mathrm{margin}}$ run reports a gap 27.8$\times$ that of the $A_{\mathrm{heur}}$ run in the geometric mean (29.0$\times$ in the median), of which 36.6$\times$ is the denominator and 0.76$\times$ is the solution --- which is \textit{better}, on 59\% of pairs. The reported gap moves because the denominator moves, with the solution pulling the other way (all sizes: Supplementary Table~S13).


\textit{Two success rules disagree, and the disagreement belongs to the weight.} Under the offset-normalized rule $g_{\mathrm{off}}$ $\leq$ 0.1\%, which is the line this implementation's own tables drew, the $A_{\mathrm{heur}}$ runs succeed at 81\%, 85\% and 82\% (N = 12, 16, 20); under the feasible-range rule $g_{F}$ $\leq$ 1\% they succeed at 78\%, 40\% and 6\%. The share of runs that pass the first and fail the second is 3\%, 46\% and 76\% at $A_{\mathrm{heur}}$, 0\% at every other setting at N = 16, and $\leq$ 5\% at N = 20. Read at four gap thresholds and three quality thresholds (Supplementary Table~S16), the disagreement is monotone in the gap threshold, from 0.06 to 0.89 at N = 20 and $A_{\mathrm{heur}}$ over 0.01\% to 0.2\%, and it stays between 0.03 and 0.07 at $A_{\mathrm{margin}}$ throughout. The 0.1\% line is not a threshold at which anything changes character.


\textit{The same solutions, three denominators (E0; Fig.~\ref{fig:3}).} Fig.~\ref{fig:3} (left) takes one fixed feasible portfolio of the reference N = 20 instance, the one of median rank, 2,423 of 4,845, and scores it as only A changes. Its feasible-range gap $g_{F}$ is 37.8\% at every A, because neither f(x) nor $\Delta_{F}$ depends on A. Its offset-normalized gap falls from 33.6\% at 1.1$\cdot$$A_{\mathrm{crit}}$ to 3.9\% at 10$\cdot$$A_{\mathrm{crit}}$ and to 0.22\% at $A_{\mathrm{heur}}$, and the all-states approximation gap 1 $-$ AR of \cite{ref22} falls from 1.36\% to 0.014\% over the same range. Nothing about the portfolio changed. Fig.~\ref{fig:3} (middle) is the deflation D(A) = $|$f* $-$ A$\cdot$K$^{2}$$|$/$\Delta_{F}$ for every size: below 1 near $A_{\mathrm{crit}}$ at N $\leq$ 12, linear in A, and 57$\times$, 81$\times$, 121$\times$, 174$\times$, 207$\times$, 240$\times$ and 279$\times$ at the $A_{\mathrm{heur}}$ of N = 8 to 30 (Supplementary Table~S3). The mechanism is the one that Sec.~III states. For feasible x the numerator f $-$ f* is independent of A while the denominator $|$f* $-$ A$\cdot$K$^{2}$$|$ is not, so two studies using the same formulation with different penalty weights report gaps that differ by roughly the ratio of their weights for identical portfolios. With A = N + 1 and K = $\lfloor$N/5$\rfloor$ the constant grows as $\Theta$(N$^{3}$), while the objective's measured spread over F rises only from 1.27 at N = 16 to 4.00 at N = 30, against A$\cdot$K$^{2}$ rising from 153 to 1,116. The same solution quality reads as a smaller number at larger N.


\textit{Over 30 instances (Fig.~\ref{fig:3}, right; Supplementary Table~S9).} The median deflation at $A_{\mathrm{heur}}$ rises monotonically with N, from 68$\times$ at N = 12 to 282$\times$ at N = 30, and no instance at any size escapes it: the smallest value observed anywhere is 36$\times$ and the largest 411$\times$. The heuristic weight sits a median 103--179$\times$ above the threshold at every size, so the overshoot is a property of the rule rather than of one covariance matrix. The threshold is the exact $A_{\mathrm{crit}}$ at N $\leq$ 20, over the 21 N = 12 instances with $A_{\mathrm{crit}}$ $>$ 0, and the upper bound $A_{\mathrm{safe}}$ at N $>$ 20, where $A_{\mathrm{heur}}$/$A_{\mathrm{safe}}$ is a lower bound. On the reference instance the overshoot is 148--1213$\times$, because its $A_{\mathrm{crit}}$ is unusually small at N $\leq$ 12 and the overshoot is a ratio against a threshold that can be near zero. That instance is near the median in deflation and below it in the Sharpe spread of its 0.1\% band (next paragraph), so the single-instance results understate the effect rather than exaggerate it.


\textit{Feasible solutions admitted by an offset-normalized threshold.} Exact enumeration makes the set that a reported gap admits countable. On the reference instances, a 0.1\% offset-normalized band holds 5.4\% of the feasible set at N = 16, 7.6\% at N = 20 and 24.7\% at N = 30, and the band is a step function of the threshold: at N = 30 it holds 1.5\% at 0.05\% and 94.7\% at 0.2\% (Supplementary Tables~S4--S5). The solutions it admits are not near-degenerate in any penalty-independent sense, since at a deflation of 279$\times$ the band reaches 28\% of $\Delta_{F}$ from the optimum, and they are admitted because the threshold was set on a compressed scale. The single-seed QAOA solutions at N = 20--30, reported at 0.024--0.053\%, sit at 5.7--9.2\% of $\Delta_{F}$ and at ranks 43 to 844, and a fixed universe-wide normalization widens the band rather than narrowing it (Supplementary Table~S6).


The domain metric is an auxiliary consequence rather than a solver error, since the Sharpe ratio is not what the QUBO optimizes (Sec.~III-A). Similar objective values did not guarantee similar Sharpe ratios. The N = 30 band spans 0.172 to 1.695; the single-seed N = 24--30 solutions sit 15--22\% below that of the optimum at gaps of 0.024--0.031\%; and of the 47 sub-optimal solutions of the seed-42 sweep at N $\leq$ 20, 17 are \textit{above} that of the optimum while the worst deviations reach $-$61\% (Supplementary Sec. S-IV, Fig.~S7). The deviations are of both signs, and the reported gap does not carry the sign.


\textit{The prescribed all-states ratio.} Applied literally to this QUBO, the range-normalized ratio of \cite{ref22} removes the offset but inherits the scale. $C_{\mathrm{max}}$ over all 2$^{N}$ assignments is attained by a violating assignment and grows with A, while the spread over F does not, so every feasible solution's ratio tends to 1. At $A_{\mathrm{heur}}$ the QAOA solution and the \textit{worst} feasible portfolio score 1.000000 and 0.999560 at N = 16 (Supplementary Table~S8). \cite{ref22} calibrates the metric on unconstrained instances and does not say which range a constrained problem should use.


\textit{The pattern does not need a quantum sampler.} Simulated annealing on the \textit{feasible set} (cardinality-preserving swaps, five seeds $\times$ 30 instances $\times$ 6 sizes, 900 runs) reaches the exact optimum in 895 of 900 runs at a median 3.4 s, so it leaves the metric nothing to hide. The same annealer on the penalty-encoded Q with single-bit flips, five seeds at N $\in$ \{12, 16, 20\}, reproduces the QAOA arm's pattern: every run feasible, offset-normalized gaps of 0.00--0.59\%, Sharpe ratios up to 64\% below that of the optimum, and ranks as deep as 4,307 of 4,845 (Supplementary Table~S11). What produces the pattern is the combination of the encoding's offset with an offset-sensitive score, and a classical solver scored the same way inherits it. The two arms are not on matched budgets (50 ADAM steps against 50,000 annealing moves), so this is a check on whether the reporting critique needs an imperfect solver, not a verdict on which method is better. The seed-42 QAOA arm over the same 30 instances reaches the exact optimum in 26, 17 and 0 of 30 at N = 12, 16 and 20 (Supplementary Table~S12), and the configuration measured here trails the classical baseline on quality as on runtime.


\subsection{Another optimizer setting, another window (E3, E5)}


The baseline runs ADAM at a fixed step of 0.1 for 50 steps. One objection to Sec.~V-A through Sec.~V-C is that a better optimizer would change the picture. The single-seed depth sweep and the steps $\times$ p $\times$ seed grid (Supplementary Sec. S-VI) had already shown three things: that p = 2 is where depth stops helping monotonically, that $\langle$$H_{\mathrm{norm}}$$\rangle$ saturates by 50--100 ADAM steps and does not improve with an 8$\times$ larger budget, and that a run can barely optimize and still report a small best-of-1,000 gap ($\langle$$H_{\mathrm{norm}}$$\rangle$ = $-$0.91 against $-$5.01 in the same cell, gaps 0.038\% and 0.022\%). Which of initialization, local minima and expressivity produces the non-monotonicity in p is not identified by those runs. E3 asks the narrower question for p = 2: on a block fixed before it was run, do two other optimizer settings change the acquisition quantities? The levels, the paired differences with their intervals, the budgets and the optimum hits are in Supplementary Table~S22, and the text below reports what they show.


The controls did not yield a uniform improvement in acquisition probability. \textit{L-BFGS-B} reduced the optimized expectation modestly and consistently, below that of the baseline in 59 of 60 instance-cells and at 1.2--1.6$\times$ the wall-clock. Its effect on the batch probability $Q_{\tau}$(1000) is of both signs and not uniformly small, from +0.009 [+0.004, +0.009] at N = 12 and $A_{\mathrm{heur}}$ to $-$0.016 [$-$0.027, +0.002] at N = 16 and $A_{\mathrm{margin}}$ (per-cell levels, intervals and budgets: Supplementary Tables~S22, S23). Its paired median changes to the one-shot $P_{\tau}$ are within $\pm$4$\cdot$10$^{-4}$ in every cell. Q compounds P over a thousand draws ($\partial$$Q_{\tau}$(S)/$\partial$$P_{\tau}$ = S(1 $-$ $P_{\tau}$)$^{S-1}$), so a change in $P_{\tau}$ of 4$\cdot$10$^{-4}$ can move $Q_{\tau}$(1000) by a tenth, and a paper that reports only $P_{\tau}$ has not reported the quantity that its own selection rule uses. ``Stopped'' is not ``converged'': the block ran at single precision, and a diagnostic re-run of a fixed sub-block at both precisions shows every single-precision run that stopped on a tolerance stopping on the \textit{objective} criterion, several of them after a failed line search, against 11 of 24 on the gradient criterion in double precision. What the precisions agree on, and the two pairs that reach different optimization outcomes from the same initialization, are in Supplementary Table~S36 and Sec. S-VI. In both of those pairs $P_{\tau}$ stays below the instance's uniform-F probability at either precision, so the comparison against uniform feasible sampling keeps its direction.


\textit{ADAM at step 0.03 with twice the budget} did not improve on the baseline within the tested budget at N = 20. Its expectation after 100 steps is higher than the baseline's after 50 in 10 of 10 instances at both weights, by 0.32 at $A_{\mathrm{margin}}$. Its $Q_{\tau}$(1000) is lower by 0.006 and by 0.106 [$-$0.187, $-$0.060], worse on 10 of 10 instances in both cells, which is the largest effect any control has on any quantity in this block, and it is an optimizer setting. Whether those runs were still descending or had settled in a different local minimum is not identified by the trajectories, which record no stopping reason for ADAM.


\textit{The scale of the seed.} The baseline's own spread of $\langle$$H_{\mathrm{norm}}$$\rangle$ over the five seeds \textit{within an instance} has a median of 3.4--5.5 at every N and both weights, and moves $Q_{\tau}$(1000) by 0.2--0.9 (Sec.~V-A). At this budget and p = 2, a large within-instance variation across initial angles thus remained, alongside optimizer-dependent changes in acquisition probability. What is compared are observed ranges: a paired median against a within-instance max $-$ min. It is not a ranking of factor contributions. Where a control moves a run by order 1, as ADAM at 0.03 does at N = 20, it has reached a different optimization outcome, in either direction. Whether a different initialization rule \cite{ref32}, \cite{ref33} or a larger budget changes the picture is not tested.


\textit{A second market period (E5; Supplementary Table~S15).} Everything above is conditional on one three-year window. E5 re-ran the pre-fixed 300-run block on the adjacent window, where $\Delta_{F}$ is 23--34\% wider at every size and one N = 12 instance whose $A_{\mathrm{crit}}$ is zero in the paper's window has a positive threshold. The main contrast, the disagreement of the two success rules, and the uniform-F comparison keep their direction (Supplementary Table~S15; the conditional-distribution and initial $\rightarrow$ final analyses were not repeated on this window). At N = 20, $A_{\mathrm{margin}}$ against $A_{\mathrm{heur}}$ raises $Q_{\tau}$(1000) on 10 of 10 instances in both windows (+0.09 and +0.12), lowers $g_{F}$ on 9 and on 7, and worsens $g_{\mathrm{off}}$ on 9 and on 10, with the denominator carrying the sign of the $g_{\mathrm{off}}$ ratio in both. The two success rules disagree in both windows: 74\% and 66\% of N = 20 $A_{\mathrm{heur}}$ runs pass $g_{\mathrm{off}}$ $\leq$ 0.1\% and fail $g_{F}$ $\leq$ 1\%. The optimized state beats the uniform-over-F sampler in 0 of 50 $A_{\mathrm{heur}}$ runs per size in the earlier window, as it does in 0 of 150 in the paper's. The levels move, which is what a different instance set produces; the signs and orderings do not. This is one further window from the same universe and asset class rather than a sample of market periods (Sec.~VII), and the full comparison is Supplementary Table~S15.


\subsection{Numerical agreement and sampling variability (E4-A, hardware)}


\textit{Numerics or sampling (E4-A; Supplementary Fig.~S14, Table~S30).} Single runs of the same configuration on the local GPU and on SV1 return bitwise-identical portfolios at N $\leq$ 12 and different feasible portfolios at N = 16 and N = 20. Three submissions of one configuration to SV1 return three different portfolios at N $\geq$ 16, and matching the precision does not restore agreement (Supplementary Sec. S-VII). Those single runs cannot say whether the output distribution differs between implementations or the same distribution was sampled differently. E4-A fixes the angles and asks both questions separately. What is compared is the exact probability of each computational basis state rather than the quantum state, since equal basis probabilities leave the relative phases unconstrained and probabilities are what the selection rule consumes. The comparison is on F $\cup$ \{$\perp$\}: the feasible bitstrings together with one symbol for ``this shot is infeasible'', $p_{\perp}$ = 1 $-$ $P_{F}$. That $p_{\perp}$ is a one-shot quantity. The probability that a whole batch of 1,000 shots holds no feasible solution is (1 $-$ $P_{F}$)$^{1000}$, and the two are not the same event.


\textit{The double-precision implementations agree; single precision differs by 10$^{-7}$.} On F $\cup$ \{$\perp$\} the managed cloud simulator's exact probabilities differ from the GPU double-precision reference by at most 2.6$\cdot$10$^{-14}$ in total variation at every one of the 12 circuits, and those of the single-threaded CPU by at most 4.3$\cdot$10$^{-15}$. Single precision, the precision of every E1 run, is the only arm that moves the distribution, by up to 6.1$\cdot$10$^{-7}$ (per-state and relative distances: Supplementary Table~S30). How much such a distance can change the distribution of the \textit{selected} output is bounded directly: for one-shot distributions at total variation d and a deterministic selection rule, a maximal coupling of the S shots puts the selected outputs within 1 $-$ (1 $-$ d)$^{S}$ $\leq$ S$\cdot$d (derivation in Supplementary Sec. S-VII). At S = 1,000 the selected-output distributions differ by at most 2.6$\cdot$10$^{-11}$ for SV1 and 6.1$\cdot$10$^{-4}$ for the single-precision arm, relative to the double-precision GPU reference. These bounds quantify how far the arithmetic can move the distribution of the returned portfolio at fixed angles. They do not bound how often two \textit{independently} sampled batches return different portfolios, which is a different quantity: two identical fair coins are at d = 0 and still disagree on half of independent draws. That is what the next paragraph measures.


\textit{Repeated batches return different portfolios.} Ten batches of 1,000 shots from one state on one backend return 9 or 10 distinct best-feasible portfolios in 10 batches at N = 20, on every arm and in every circuit: 10 in 15 of the 16 arm $\times$ circuit cells, and 9 feasible portfolios in the one SV1 cell where a batch held no feasible shot (10 selected outcomes there if $\perp$ is counted). At the smaller sizes the count depends on the state as much as on the size: at N = 12 it is 1--2 for the seed-42 states and 5--10 for the seed-43 states, and at N = 16 it is 3--5 and 1--5 (Supplementary Table~S30). In no N = 20 batch is the best feasible shot the exact optimum. The spread is what i.i.d. sampling from the arm's own exact distribution predicts. A Monte Carlo of 10 batches from each arm's distribution puts the observed number of distinct portfolios inside its 95\% interval in 46 of 48 arm $\times$ circuit cells, and the feasible-shot count of 456 of 480 batches lies within two binomial standard deviations of 1,000$\cdot$$P_{F}$.


Repeated sampling from a fixed distribution was therefore sufficient to produce portfolio disagreement of the magnitude that the single runs showed, while the arithmetic can move the selected-output distribution by at most the bounds above. No individual disagreement is attributed to either channel. E4-A does not test the third possibility, which is different final angles from optimization runs whose trajectories differ by $\sim$10$^{-6}$ per step on the two backends, so sampling is shown to be sufficient rather than sole. The consequence for a device-abstraction layer is that backend-independence of the QUBO does not imply backend-independence of the returned portfolio, and that agreement at small sizes does not establish it at larger ones, where the feasible set is large.


\textit{Hardware: four tasks on IonQ Forte-1.} Each task samples 100 shots from a circuit whose angles were optimized on the local GPU (two seeds at N = 8, one at N = 12, one at N = 16; \$33.20 in all). In the N = 12 task, 89 of 100 shots violated the constraint and the remaining 11 contained the exact optimum, so a best-of-shots report would record a 0.0000\% gap. At N = 16, 18 shots were feasible and the best was 0.0438\% away. The two N = 8 tasks, differing only in initial angles, returned 25\% and 7\% feasible shots. The four tasks illustrate the distinction between the quality of the best feasible outcome and the frequency of feasible outcomes. They identify no size limit and no ranking of noise against instance structure, which four tasks at one calibration could not, and the product-fidelity model that predicted a ceiling near N = 12 was contradicted by the N = 16 task (Supplementary Sec. S-VIII).


\subsection{Resources and how the measured quantities scale}


\textit{Scale and reference.} The single-instance, single-seed benchmark (N = 8--30, seed 42; Supplementary Table~S2) locates the simulator envelope: classical runtime is flat in the seconds range, while QAOA grows from 30 s at N = 20 to 3.14 h at N = 30. At N = 30, single precision uses 41.3 of 49 GB of VRAM, double precision does not fit, and N = 50 ($\approx$ 9 PB for a dense statevector) is beyond any memory this study could reach. Its gaps are 0.000\% at N $\leq$ 16 and 0.024--0.053\% at N = 20--30, with every row feasible, and Sec.~V-D measures what those numbers denominate.


\textit{How the measured quantities scale.} Over the sizes the cross covers, the quantities of Sec.~III-C fall with N at about the rate at which the constraint-respecting control falls. The median one-shot success probability at $A_{\mathrm{heur}}$ drops 43-fold from N = 12 to N = 20 and that of the uniform-feasible draw drops 37-fold, while the feasible set grows 73-fold (66, 560 and 4,845). The circuit's resource count grows as N qubits and p$\cdot$[C(N, 2) + 2N] parameterized rotations, 460 at N = 20 with p = 2, and the evaluation's own cost grows with the feasible set, since the exact optimum, the exact rank and $\Delta_{F}$ come from enumerating C(N, K) candidates (5.9$\cdot$10$^{5}$ of them in 5.7 s at N = 30). Both the rotation count and the enumeration are reported as what this study measured, not as native gate counts, as fault-tolerant resource estimates, or as an extrapolation to sizes where the feasible set cannot be enumerated (Sec.~VI-A, Sec.~VII).


\section{Engineering Implications}


The measurements identify three engineering requirements: formulation-independent scoring, feasibility-aware acquisition metrics, and a recorded execution configuration. Runtime comparisons give context for the evaluated implementation rather than evidence about quantum advantage: on a statevector simulator, the exact classical baseline dominated on runtime by three to six orders of magnitude, and on solution quality at every size we could run (Sec.~V-D, Sec.~V-G). What the measurement leaves behind is reusable: what to report, what the same experiment cost on metered backends, the implementation defects that changed results before they were caught, and an artifact that lets a reader re-run any of it (Sec.~IV-D and the availability statement). The infrastructure, its cost guards, and the provider-specific defects met in deploying it are documented in the artifact rather than here.


\subsection{What to report}


Three reporting proposals exist \cite{ref22}, \cite{ref23}, \cite{ref24}, and none of them is a formal standard. We adopt the fields of \cite{ref22}, \cite{ref23} where they apply; Table~\ref{tab:V} collects what the measurements here showed a report has to carry, as recommendations rather than as a standard. Table~\ref{tab:V} lists what each result of Sec.~V showed to be necessary, linking each reporting item to the measurement that motivates it.


Which of those items survive where the feasible set cannot be enumerated is the next question, since every quantity above that names f*, $\Delta_{F}$ or a rank in F is computed here by enumeration. The feasible fraction of a batch, the original objective of the returned solution, the task counts and the cost, and a comparison against a feasible classical baseline need none of it. A success probability can still be estimated against a target fixed in advance rather than against the optimum, with its sampling uncertainty reported; what cannot be had are $g_{F}$, the rank, and $D_{\mathrm{cond}}$, which are defined against the whole feasible set. An estimate of the first kind should be labeled as such, and not read as one of the exact quantities this paper reports.





\begin{table*}[!t]
\centering
\scriptsize
\caption{Reporting items, and the measurement that motivates each.}
\label{tab:V}
\begin{tabular}{p{0.056\textwidth}p{0.337\textwidth}p{0.517\textwidth}}
\toprule
Evaluated quantity & Report & Result that showed the need \\
\midrule
Solution quality & $g_{F}$ against the objective's spread over F, or the rank in F; the penalty weight A alongside any offset-normalized figure & E0: the same portfolio reads 33.6\% or 0.22\% as A alone changes; the all-states ratio of \cite{ref22} is non-discriminative at $A_{\mathrm{heur}}$ (Sec.~V-D) \\
Optimization & the optimized expectation beside the best-of-shots figure (\cite{ref22}'s mean/best split) & a run at $\langle$$H_{\mathrm{norm}}$$\rangle$ = $-$0.91 and one at $-$5.01 both report small offset-normalized gaps, 0.038\% and 0.022\% (Sec.~V-E) \\
Acquisition & $P_{F}$, $P_{\tau}$ and $Q_{\tau}$(S) at the budget used, against a constraint-respecting control sampler & E2: optimization raises the median $Q_{\tau}$(1000) from 0.06 to 0.34 at N = 16, yet at $A_{\mathrm{heur}}$ and $\tau$ = 1\% beats a uniform feasible draw's $P_{\tau}$ in 0 of the 450 runs over the three sizes (Sec.~V-B) \\
Execution conditions & the initial-angle seed distribution, precision, backend and sampler seeding, as factors & one seed reverses the main contrast (Sec.~IV-E); single runs on two backends disagree at N $\geq$ 16 (Sec.~V-F) \\
Cost & measured task counts and billed time per run & one sampling call was billed as 177 tasks, not one, until its angles were marked non-trainable (Sec.~VI-B) \\
Domain metric & the finance quantity as an auxiliary, with its definition & similar objective values do not guarantee similar Sharpe ratios (Sec.~V-D) \\
\bottomrule
\end{tabular}
\end{table*}


The sharpest form of the first three rows is that the penalty weight is a free parameter of the reported metric (Table~\ref{tab:IV}). Re-scoring fixed feasible solutions at a larger penalty weight reduces the offset-normalized gap without changing the solutions or their sampling distribution (E0). In the re-optimization contrast, the larger weight also yielded a lower acquisition probability under the tested settings (E1). Those are two effects, one of the number and one of the run, and neither is a law that a larger A always degrades the sampler. Reporting A is the minimum. Reporting against the penalty-free objective, or reporting a rank, removes the degree of freedom, and reporting $Q_{\tau}$(S) against the uniform-feasible control says whether the optimizer did anything that a constraint-aware guess would not. The case of Mancilla et al. \cite{ref27} makes the cost of under-reporting concrete. A realized Sharpe of 1.81 for QAOA against 1.31 for annealing is read there as a property of the energy landscape, with no objective value for a single selected portfolio and no exact optimum --- and C(10, 5) = 252 candidates is small enough to enumerate directly.


\subsection{Compute and billing}


Table~\ref{tab:II} reports, for every metered run, the task count and the dollars billed, and the totals of the study (\$37.90 on Braket). Task counts are reported because one of them was wrong by two orders of magnitude before it was measured: in the N = 8 SV1 pilot, the final 1,000-shot sampling call was billed as 177 tasks rather than one. The optimizer's angles reached that call still marked trainable, so the framework returned 176 shifted tapes for a gradient that no one reads (2$\cdot$p$\cdot$[C(N, 2) + 2N] at N = 8, reproduced locally with no device; 920 at N = 20). The run's 227 tasks became 51 once the angles were frozen. The mechanism is inferred from the count and the fix rather than from an instrumented re-run, and the per-run ledger is in the artifact. On a metered backend, the executed circuit count is not implied by the algorithm's description, so it has to be measured and reported per run.


\subsection{Checks on the recorded results}


The results of Sec.~V rest on four correspondences that were checked rather than assumed, each of them because an implementation defect had at some point broken it. (i) \textit{Objective against QUBO.} The evaluator of Sec.~III-C scores from the returns and covariance, not from the QUBO matrix, and it agrees with the QUBO path on every state at N = 8 and 12 (E0). (ii) \textit{State against batch.} The analytic state quantities and the final sampling call receive the same final parameters, which regression tests pin on both optimizer paths. Every arm is reported with two consistency checks: the batch's feasible fraction must track the final state's $P_{F}$, and the batch counts must fall within what the recorded states predict under independent draws (Supplementary Table~S22). (iii) \textit{Sampler against seed.} E4-A seeds the device's own generator, because the Lightning devices draw a seed from NumPy's global generator once, at construction. (iv) \textit{Hardware records against device schema.} The four IonQ task records are read from the device's nested result schema (Supplementary Sec. S-VIII).


One of the seven defects reached a result reported in Sec.~V: the E3 L-BFGS-B arm's sampler read the initial parameter vector, so its batch columns described unoptimized states. That arm was re-run with the fix on the same host. Every state of the 300 runs reproduced bit for bit, and the re-run's counts sit within 3.0 hits and 2.7 no-feasible runs of the prediction in every cell. The defects in the order found, what each would have reported, and the checks that caught them are in Supplementary Sec. S-IX.


\section{Limitations and Threats to Validity}


\textbf{Formulation and data.} Every measurement is on one constrained problem, K of N by cardinality, under two encodings of it: the quadratic penalty at six weights, and the constraint-preserving ansatz of E6 at A = 0. Slack variables or a Lagrangian update would place the constant differently or remove it, and would change $P_{F}$ by construction; an implementation that restores the constant and scores on the original objective would not report $g_{\mathrm{off}}$ at all. The cross covers N $\leq$ 20, where the feasible set is enumerable. The census extends to N = 30 by enumeration on 30 instances per size (Supplementary Table~S9), the QAOA runs above N = 20 are single, and N = 50 was not enumerated in this study. All E0--E4 instances are subsets of one 50-ticker Nikkei universe over one three-year window, so every interval in Sec.~V is conditional on that universe, that window, and that drawing rule (Sec.~IV-E). E5 adds one adjacent window of the same universe, chosen by adjacency rather than sampled, and it does not make the intervals unconditional. Thirty subsets is fewer than the 100 instances of Uotila et al. \cite{ref18} or the 250 of Stopfer and Wagner \cite{ref13}. The Sharpe ratio is auxiliary: it is measured on the raw series for the equal-weight portfolio, it is not what the QUBO optimizes, and its deviations from those of the optimum are reported as domain-level consequences rather than as predictions.


\textbf{Optimization.} Five seeds per instance show that a within-instance tail exists; they do not estimate how often it occurs. The design was not sized from a minimum detectable effect, and the intervals are conditional on the fixed seed set, one member of which reverses the main contrast on its own at each of the two larger sizes (Supplementary Table~S19). E3 tests two settings, not the optimizer space. One of them did not improve on the baseline within its budget, and the other's single-precision terminations include line-search failures (Sec.~V-E). A schedule, a larger budget, a different initialization rule \cite{ref32}, \cite{ref33}, or depth beyond 2 with a warm start could change the relative sizes of the effects that Sec.~V-E summarizes, and the depth non-monotonicity remains an observation whose cause is not identified. The multistart comparison is at equal shots, with the optimization cost stated beside it, rather than at equal total cost. $A_{\mathrm{crit}}$ and $A_{\mathrm{margin}}$ are reference conditions computed by enumeration, not procedures that scale. The critical penalty is a ground-state feasibility threshold, not a guarantee of reliable finite-shot acquisition: near it, infeasible batch minima are common (55\% of N = 20 runs at 1.1$\cdot$$A_{\mathrm{crit}}$), which is what makes the feasibility-aware selection rule essential. These observations do not by themselves rule out near-threshold weights under that rule, whose utility is assessed through the acquisition probability and the solution quality at the chosen budget (Sec.~V-C). E2's strongest control is matched on candidates generated, not on cost. That the optimized state at $A_{\mathrm{heur}}$ beat a uniform draw from F in none of 450 runs at the tested thresholds says that the best-of-shots figure cannot confirm the optimization's contribution at these sizes; it does not say that a constraint-aware quantum method would fail to. $Q_{\tau}$(S) assumes i.i.d. shots.


\textbf{Backend and hardware.} E4-A separates the output distribution from the draws, not the optimization: 12 circuits at one weight, evaluated at fixed angles, bound how far the arithmetic can move the selected-output distribution, and they show sampling to be sufficient for the observed disagreement. Whether an optimization run on SV1 reaches the same angles as the GPU run does was not tested (E4-B, $\approx$1,500 tasks), so the backend attribution is partial. Hardware is four tasks at one calibration: two seeds at N = 8, one each at N = 12 and 16, 100 shots each. They establish no size limit, no noise ranking, and no distribution. Continuous weights and richer constraints are deferred.


\section{Conclusion}


The contribution of a constrained quantum optimizer reads as two factors: how often it satisfies the constraint and which of the feasible solutions it prefers. Under the dominate-the-objective rule that this implementation uses, the second was near-uniform. The median preference ratio inside the feasible set remained within 1\% of unity, where the run's initial state already was, and across 900 runs, over three optimizer settings and two price windows, the optimized state's one-shot success probability never exceeded that of a uniform draw from that set. What optimization changed was the first factor: the feasible probability $P_{F}$ rose on 30 of 30 instance medians at each size. The preference appears when the objective carries more of the cost, and the share of runs exceeding the control rises monotonically with $\epsilon_{F}$, to 41\% at 1.1 times the feasibility threshold at 20 assets. Moving the constraint into the ansatz instead puts 420 of 450 runs above the control, whereas a block that only removes the penalty gains the preference without the mass. It is the product of the two factors, not either one alone, that rises above the control.


The evidence comes from one constrained problem at N $\leq$ 20 and p = 2, under a penalty weight swept over six settings and a constraint-preserving ansatz beside it, on two adjacent windows of one 50-ticker universe, with two optimizers in three settings. Within it, the weaker weights are not a recommendation; the median run stays below the control even at the weakest, the threshold is computed by enumeration, and 55\% of batches there have an infeasible minimum at 20 assets. None of these limits bounds QAOA or argues for or against a quantum speedup, a question that these sizes cannot address. By contrast, enumeration solves instances of this size outright.


What these measurements do influence is how such a run is reported. The offset-normalized gap registers none of this and moves the other way: lowering the weight worsens it in 29 of 30 instances at 20 assets, while the acquisition probability improves in 30 of 30, because the weight sits in the gap's denominator and also sets how much of the objective the circuit resolves. In our comparison, the gap ranked best the setting whose optimized state showed no preference. A benchmark that emphasizes what its quantum part contributes therefore carries the feasible and the acquisition probability at the budget used, alongside the gap and the weight, with a constraint-respecting control (Table~\ref{tab:V}), and Sec.~VI-A marks which of those items survive where the feasible set cannot be enumerated.


\section*{Data and Code Availability}

The code, the derived input statistics of the two price windows ($\mu$ and $\Sigma$ of the 50-ticker universe, from which every instance is sliced; the price snapshots themselves are third-party data and are not redistributed), the census and its per-paper cells, every result CSV this paper reads (one row per run and one per instance, with the aggregates the tables are checked against), the figure scripts and the table-checking script are at \url{https://github.com/tobitaQ/qaoa-portfolio-penalty-benchmark}, release \texttt{v1.2}, archived at doi:10.5281/zenodo.24032941; the README indexes the files by experiment. A clone reproduces every number without network access to a market-data provider, at three levels: every figure and table from the committed CSVs (seconds, no GPU); one N = 8 or N = 12 instance end to end (minutes on a CPU); the cross from the manifest (13.9 GPU h). The per-run artifacts (exact probabilities on F under the initial and final state, angles, trajectory, shot counts) are deposited in the same archive: 2,565 files and $\approx$100 MB compressed for E1, 900 files and $\approx$34 MB for E6 and its control. They are needed only to recompute the E2 columns and $D_{\mathrm{cond}}$, which the committed CSVs already hold; any full state vector is reconstructed from the recorded angles and seed. The four IonQ task records are committed in the device's own nested schema. Supplementary Figs.~S1--S14 and Tables~S1--S38 are generated by the same scripts from the same committed files. Cloud and hardware paths require an explicit backend name and a results tag.

\begin{thebibliography}{99}
\bibitem{ref1} H. Markowitz, ``Portfolio selection,'' \textit{J. Finance}, vol. 7, no. 1, pp. 77--91, Mar. 1952, doi: 10.1111/j.1540-6261.1952.tb01525.x.

\bibitem{ref2} D. Bienstock, ``Computational study of a family of mixed-integer quadratic programming problems,'' \textit{Mathematical Programming}, vol. 74, pp. 121--140, 1996, doi: 10.1007/BF02592208.

\bibitem{ref3} A. Lucas, ``Ising formulations of many NP problems,'' \textit{Frontiers in Physics}, vol. 2, art. 5, 2014, doi: 10.3389/fphy.2014.00005.

\bibitem{ref4} E. Farhi, J. Goldstone, and S. Gutmann, ``A quantum approximate optimization algorithm,'' arXiv:1411.4028, 2014.

\bibitem{ref5} R. Or\'{u}s, S. Mugel, and E. Lizaso, ``Quantum computing for finance: Overview and prospects,'' \textit{Reviews in Physics}, vol. 4, art. 100028, 2019, doi: 10.1016/j.revip.2019.100028.

\bibitem{ref6} D. J. Egger \textit{et al.}, ``Quantum computing for finance: State-of-the-art and future prospects,'' \textit{IEEE Trans. Quantum Eng.}, vol. 1, pp. 1--24, 2020, doi: 10.1109/TQE.2020.3030314.

\bibitem{ref7} D. Herman \textit{et al.}, ``Quantum computing for finance,'' \textit{Nature Reviews Physics}, vol. 5, pp. 450--465, 2023, doi: 10.1038/s42254-023-00603-1.

\bibitem{ref8} S. Brandhofer, D. Braun, V. Dehn, G. Hellstern, M. H\"{u}ls, Y. Ji, I. Polian, A. S. Bhatia, and T. Wellens, ``Benchmarking the performance of portfolio optimization with QAOA,'' \textit{Quantum Inf. Process.}, vol. 22, no. 1, art. 25, 2023; arXiv:2207.10555.

\bibitem{ref9} J. S. Baker and S. K. Radha, ``Wasserstein solution quality and the quantum approximate optimization algorithm: A portfolio optimization case study,'' arXiv:2202.06782, 2022.

\bibitem{ref10} L. Lozano, ``A penalty-free pipeline for direct quantum-annealer portfolio optimization,'' \textit{arXiv preprint} arXiv:2605.17628, May 2026.

\bibitem{ref11} S. Hadfield, Z. Wang, B. O'Gorman, E. G. Rieffel, D. Venturelli, and R. Biswas, ``From the quantum approximate optimization algorithm to a quantum alternating operator ansatz,'' \textit{Algorithms}, vol. 12, no. 2, art. 34, 2019, doi: 10.3390/a12020034.

\bibitem{ref12} Z. Wang, N. C. Rubin, J. M. Dominy, and E. G. Rieffel, ``XY mixers: analytical and numerical results for the quantum alternating operator ansatz,'' \textit{Phys. Rev. A}, vol. 101, art. 012320, 2020, doi: 10.1103/PhysRevA.101.012320.

\bibitem{ref13} E. Stopfer and F. Wagner, ``Quantum portfolio optimization: An extensive benchmark,'' arXiv:2509.17876, 2025 (rev. Jul. 2026).

\bibitem{ref14} D. Goldberg, ``What every computer scientist should know about floating-point arithmetic,'' \textit{ACM Comput. Surv.}, vol. 23, no. 1, pp. 5--48, Mar. 1991, doi: 10.1145/103162.103163.

\bibitem{ref15} G. Rosenberg, P. Haghnegahdar, P. Goddard, P. Carr, K. Wu, and M. L\'{o}pez de Prado, ``Solving the optimal trading trajectory problem using a quantum annealer,'' \textit{IEEE J. Sel. Topics Signal Process.}, vol. 10, no. 6, pp. 1053--1060, Sep. 2016, doi: 10.1109/JSTSP.2016.2574703.

\bibitem{ref16} S. Mugel \textit{et al.}, ``Dynamic portfolio optimization with real datasets using quantum processors and quantum-inspired tensor networks,'' \textit{Phys. Rev. Research}, vol. 4, art. 013006, 2022, doi: 10.1103/PhysRevResearch.4.013006.

\bibitem{ref17} D. Venturelli and A. Kondratyev, ``Reverse quantum annealing approach to portfolio optimization problems,'' \textit{Quantum Machine Intelligence}, vol. 1, no. 1, pp. 17--30, 2019, doi: 10.1007/s42484-019-00001-w.

\bibitem{ref18} V. Uotila, J. Ripatti, and B. Zhao, ``Higher-order portfolio optimization with quantum approximate optimization algorithm,'' in \textit{Proc. IEEE Int. Conf. Quantum Computing and Engineering (QCE)}, 2025, pp. 2238--2249, doi: 10.1109/QCE65121.2025.00244.

\bibitem{ref19} A. Verma and M. Lewis, ``Penalty and partitioning techniques to improve performance of QUBO solvers,'' \textit{Discrete Optimization}, vol. 44, art. 100594, 2022, doi: 10.1016/j.disopt.2020.100594.

\bibitem{ref20} M. Ayodele, ``Penalty weights in QUBO formulations: permutation problems,'' in \textit{Proc. EvoCOP 2022}, LNCS 13222, pp. 159--174; preprint arXiv:2206.11040.

\bibitem{ref21} J. A. Montanez-Barrera, D. Willsch, A. Maldonado-Romo, and K. Michielsen, ``Unbalanced penalization: a new approach to encode inequality constraints of combinatorial problems for quantum optimization algorithms,'' \textit{Quantum Sci. Technol.}, vol. 9, no. 2, art. 025022, 2024; arXiv:2211.13914.

\bibitem{ref22} A. Abbas \textit{et al.}, ``Challenges and opportunities in quantum optimization,'' \textit{Nature Reviews Physics}, vol. 6, no. 12, pp. 718--735, Dec. 2024, doi: 10.1038/s42254-024-00770-9.

\bibitem{ref23} T. Koch \textit{et al.}, ``Quantum optimization benchmarking library --- The intractable decathlon,'' arXiv:2504.03832, 2025.

\bibitem{ref24} L. Lozano, ``Where the quantum lives in D-Wave hybrid portfolio optimization: an operational decomposition audit,'' \textit{arXiv preprint} arXiv:2605.17623, 2026.

\bibitem{ref25} M. Hodson, B. Ruck, H. Ong, D. Garvin, and S. Dulman, ``Portfolio rebalancing experiments using the Quantum Alternating Operator Ansatz,'' \textit{arXiv preprint} arXiv:1911.05296, Nov. 2019.

\bibitem{ref26} A. Oralkhan and T. Zhaxalykov, ``Investigation of hardware architecture effects on quantum algorithm performance: A comparative hardware study,'' arXiv:2601.05286, 2026.

\bibitem{ref27} J. Mancilla, T. D. Bouloumis, and F. Goguikian, ``Constrained portfolio optimization via quantum approximate optimization algorithm (QAOA) with XY-mixers and Trotterized initialization: A hybrid approach for direct indexing,'' \textit{arXiv preprint} arXiv:2602.14827, Feb. 2026.

\bibitem{ref28} N. N. Hegade, P. Chandarana, K. Paul, X. Chen, F. Albarr\'{a}n-Arriagada, and E. Solano, ``Portfolio optimization with digitized counterdiabatic quantum algorithms,'' \textit{arXiv preprint} arXiv:2112.08347, Dec. 2021.

\bibitem{ref29} B. Wu and L. Wang, ``A two-step quantum approximate optimization algorithm for portfolio optimization and risk assessment,'' \textit{Quantum Reports}, vol. 8, no. 2, art. 45, May 2026, doi: 10.3390/quantum8020045.

\bibitem{ref30} N. Innan, A. Saleem, A. Marchisio, and M. Shafique, ``Quantum portfolio optimization with expert analysis evaluation,'' \textit{arXiv preprint} arXiv:2507.20532, Jul. 2025.

\bibitem{ref31} P. K. Barkoutsos, G. Nannicini, A. Robert, I. Tavernelli, and S. Woerner, ``Improving variational quantum optimization using CVaR,'' \textit{Quantum}, vol. 4, art. 256, 2020, doi: 10.22331/q-2020-04-20-256.

\bibitem{ref32} L. Zhou, S.-T. Wang, S. Choi, H. Pichler, and M. D. Lukin, ``Quantum approximate optimization algorithm: Performance, mechanism, and implementation on near-term devices,'' \textit{Phys. Rev. X}, vol. 10, art. 021067, 2020, doi: 10.1103/PhysRevX.10.021067.

\bibitem{ref33} S. H. Sack and M. Serbyn, ``Quantum annealing initialization of the quantum approximate optimization algorithm,'' \textit{Quantum}, vol. 5, art. 491, Jul. 2021, doi: 10.22331/q-2021-07-01-491.

\bibitem{ref34} J. R. McClean, S. Boixo, V. N. Smelyanskiy, R. Babbush, and H. Neven, ``Barren plateaus in quantum neural network training landscapes,'' \textit{Nat. Commun.}, vol. 9, art. 4812, 2018, doi: 10.1038/s41467-018-07090-4.

\bibitem{ref35} L. Bittel and M. Kliesch, ``Training variational quantum algorithms is NP-hard,'' \textit{Phys. Rev. Lett.}, vol. 127, no. 12, art. 120502, 2021, doi: 10.1103/PhysRevLett.127.120502.

\bibitem{ref36} G. G. Guerreschi and A. Y. Matsuura, ``QAOA for Max-Cut requires hundreds of qubits for quantum speed-up,'' \textit{Scientific Reports}, vol. 9, art. 6903, 2019, doi: 10.1038/s41598-019-43176-9.

\bibitem{ref37} Amazon Web Services, ``Amazon Braket Developer Guide,'' AWS documentation. [Online]. Available: {\urlstyle{rm}\url{https://docs.aws.amazon.com/braket/}} (accessed Nov. 7, 2026).

\bibitem{ref38} V. Bergholm \textit{et al.}, ``PennyLane: Automatic differentiation of hybrid quantum-classical computations,'' arXiv:1811.04968, 2018.

\bibitem{ref39} H. Bayraktar \textit{et al.}, ``cuQuantum SDK: A high-performance library for accelerating quantum science,'' in \textit{Proc. IEEE Int. Conf. Quantum Computing and Engineering (QCE)}, 2023, pp. 1050--1061, doi: 10.1109/QCE57702.2023.00119.

\bibitem{ref40} T. Jones and J. Gacon, ``Efficient calculation of gradients in classical simulations of variational quantum algorithms,'' arXiv:2009.02823, 2020.

\bibitem{ref41} C. Collberg and T. A. Proebsting, ``Repeatability in computer systems research,'' \textit{Communications of the ACM}, vol. 59, no. 3, pp. 62--69, Mar. 2016, doi: 10.1145/2812803.

\bibitem{ref42} G. K. Sandve, A. Nekrutenko, J. Taylor, and E. Hovig, ``Ten simple rules for reproducible computational research,'' \textit{PLoS Computational Biology}, vol. 9, no. 10, art. e1003285, 2013, doi: 10.1371/journal.pcbi.1003285.

\bibitem{ref43} Association for Computing Machinery, ``Artifact review and badging --- version 1.1,'' ACM policy, Aug. 2020. [Online]. Available: {\urlstyle{rm}\url{https://www.acm.org/publications/policies/artifact-review-and-badging-current}} (accessed Nov. 7, 2026).

\bibitem{ref44} S. Kirkpatrick, C. D. Gelatt, and M. P. Vecchi, ``Optimization by simulated annealing,'' \textit{Science}, vol. 220, no. 4598, pp. 671--680, 1983, doi: 10.1126/science.220.4598.671.

\bibitem{ref45} D. P. Kingma and J. Ba, ``Adam: A method for stochastic optimization,'' in \textit{Proc. 3rd Int. Conf. Learning Representations (ICLR)}, 2015. arXiv:1412.6980.

\bibitem{ref46} R. H. Byrd, P. Lu, J. Nocedal, and C. Zhu, ``A limited memory algorithm for bound constrained optimization,'' \textit{SIAM J. Sci. Comput.}, vol. 16, no. 5, pp. 1190--1208, 1995, doi: 10.1137/0916069.

\bibitem{ref47} P. Virtanen et al., ``SciPy 1.0: fundamental algorithms for scientific computing in Python,'' \textit{Nature Methods}, vol. 17, pp. 261--272, 2020, doi: 10.1038/s41592-019-0686-2.

\bibitem{ref48} A. B\"{a}rtschi and S. Eidenbenz, ``Deterministic preparation of Dicke states,'' in \textit{Fundam. Comput. Theory}, LNCS 11651. Springer, 2019, pp. 126--139, doi: 10.1007/978-3-030-25027-0\_9.



\end{thebibliography}

\EOD

\end{document}
